\documentclass[a4paper,12pt]{report}
\setcounter{secnumdepth}{5}
\setcounter{tocdepth}{3}
\input{LaTeX-ToolKit/template.tex}
\begin{document}
\title{Linear Algebra and Functional Analysis}
\author{沈威宇}
\date{\temtoday}
\titletocdoc
\chapter{Linear Algebra (線性代數) and Functional Analysis (泛函分析)}
\section{Vector Space and Functional}
\subsection{General}
\subsubsection{Vector space (向量空間) or Linear space (線性空間)}
A vector space over a field $F$ is a non-empty set $V$ together with a binary operation and a binary function that satisfy the eight axioms listed below. In this context, the elements of $V$ are commonly called vectors, and the elements of $F$ are called scalars.

The binary operation, called vector addition or simply addition assigns to any two vectors $v$ and $w$ in $V$ a third vector in $V$ which is commonly written as $v + w$, and called the sum of these two vectors.

The binary function, called scalar multiplication, assigns to any scalar $a$ in $F$ and any vector $v$ in $V$ another vector in $V$, which is denoted $av$.

To have a vector space, the eight following axioms must be satisfied for every $u$, $v$ and $w$ in $V$, and $a$ and $b$ in $F$.
\begin{enumerate}
\item Associativity of vector addition: $u + (v + w) = (u + v) + w$.
\item Commutativity of vector addition: $u + v = v + u$.
\item Identity element of vector addition: There exists an element $ 0 \in V$, called the zero vector, such that $\forall v\in V\colon v+ 0=v$.
\item Inverse elements of vector addition: $\forall v\in V\colon \exists$ an element $-v\in V$, called the additive inverse or inverse of $v$, such that $v + (−v) = 0$.
\item Compatibility of scalar multiplication with field multiplication: $a(bv) = (ab)v$.
\item Identity element of scalar multiplication: $1v = v$, where 1 denotes the multiplicative identity in $F$.
\item Distributivity of scalar multiplication with respect to vector addition: $a(u + v) = au + av$.
\item Distributivity of scalar multiplication with respect to field addition: $(a + b)v = av + bv$.
\end{enumerate}
When the scalar field is the real numbers, the vector space is called a real vector space, and when the scalar field is the complex numbers, the vector space is called a complex vector space. These two cases are the most common ones, but vector spaces with scalars in an arbitrary field $F$ are also commonly considered. Such a vector space is called an $F$-vector space or a vector space over $F$.
\subsubsection{Linear combination (線性組合)}
A linear combination of a series of vectors $(v_1,\ldots,v_n)$ in a vector space $V$ over a field $F$ is any vector of the form
\[\sum_{i=1}^nc_iv_i\]
where $(c_1,\ldots,c_n)$ is a series of scalars in $F$ called coefficients.
\subsubsection{Isomorphism (同構)}
Vector spaces $X$ and $Y$ are called isomorphic if there exists a mapping $f\colon X\to Y$ between them such that $f$ is bijective and continuous and $f^{-1}$ is continuous, written as \( X \cong Y \), and $f$ is called a isomorphism between them.
\subsubsection{Division}
For nonzero vector $a$ and vector $b$ in a vector space such that for some scalar $s$, $sa=b$, $\frac{b}{a}\coloneq s$.
\subsubsection{Affine space (仿射空間)}
An affine space is a point set $A$ together with a vector space $\ora{A}$, and a transitive and free action of the additive group of $\ora{A}$ on the set $A$. The vector space $\ora{A}$ is said to be associated to the affine space, and its elements are called vectors, translations, or sometimes free vectors.

Explicitly, the definition above means that the action is a mapping, generally denoted as an addition,
\[\begin{aligned}
& A\times\ora{A}\to A\\
& (a,v)\mapsto a+v
\end{aligned}\]
that has the following properties.
\begin{enumerate}
\item Right identity: \[\forall a\in A\colon a+0=a,\quad\tx{where $ 0$ is the zero vector in $\ora{A}$}\]
\item Associativity: \[\forall v,w\in \ora{A},\forall a\in A,\;(a+v)+w=a+(v+w)\]
\item Free and transitive action: \[\forall a\in A\colon \tx{the mapping }\ora{A}\to A\colon v\mapsto a+v\tx{ is a bijection.}\]
\item Existence of one-to-one translations \[\forall v\in \ora{A} \colon \tx{the mapping }\ora{A}\to A\colon v\mapsto a+v\tx{ is a bijection.}\]
\end{enumerate}
\subsubsection{Linear map (線性映射)/operator/transformation}
A map $T\colon V\to W$ between vector spaces $V$ and $W$ over field $F$ is linear iff for any $u,v\in V$ and scalar $c$, the following two conditions are satisfied:
\bit
\item Preservation of vector addition:
  \[T(u+v)=T(u)+T(v)\]
\item Preservation of scalar multiplication:
  \[T(cu)=cT(u)\]
\eit
The space of linear operators from $V$ to $W$ is a vector space, denoted as $L(V,W)$.
\sssc{Properties of linear maps}
For any linear map $T\colon V\to W$,
\[T(0)=0\]
\bpr
\[T(0)=T(0+0)=2T(0)\]
\epr
For any linear map $T\colon V\to W$, $u,v\in V$, and scalars $a,b$,
\[T(au+bv)=aT(u)+bT(v)\]
\bpr
\[T(au+bv)=T(au)+T(bv)=aT(u)+bT(v)\]
\epr
\subsubsection{Invertible map (可逆映射)}
A map $T\colon V\to W$ is Invertible if and only if it is bijective. The unique inverse of $T$, called $T^{-1}$, is defined to be:
\[T^{-1}\colon W\to V;\,\forall v\in V\colon T(v)\mapsto v=T^{-1}(T(v)).\]
\subsubsection{Linear subspace, vector subspace, or subspace}
A linear subspace of a vector space is a subset of it that is closed under vector addition and scalar multiplication, and thus a vector space.
\subsubsection{Linear span, linear hull, or span}
A span of a subset $S$, denoted as $\span S$ or $\langle S\rangle$, of a vector space is the smallest linear subspace containing $S$, and it is called that $S$ is a generating/spanning set of $\span S$ and $S$ generates/spans $\span S$.
\subsubsection{In the span iff new span unchanged}
Let $S$ be a set of vectors and $v$ be a vector. $v\in\span S$ iff $\span S=\span(S\cup\{v\})$.
\bpr
Only if: Suppose $v\in\span S$. We can write $v$ as a linear combination of vectors in $S$. For any $w\in\span(S\cup\{v\})$, we can write $w$ as a linear combination of vectors in $S\cup\{S\}$. By replacing $v$ with the linear combination of vectors in $S$, we can write $w$ as a linear combination of vectors in $S$.

If: Suppose $v\notin\span S$. Since $v=1v$, $v\in\span(S\cup\{v\})\setminus S$, thus $\span S\ne\span(S\cup\{v\})$.
\epr
\subsubsection{Linear independence (線性獨立)}
A set of vectors \(\{v_1,\ldots,v_n\}\) is linearly independent iff the equation:
\[\sum_{i=1}^nc_iv_i=0 \]
is true only if all coefficients \(c_1,c_2, \dots,c_n\) are zero, then the set of vectors is linearly independent; otherwise the set of vectors is linearly dependent.
\subsubsection{Linearly dependent set condition}
A set $S$ of vectors is linearly dependent iff either $\mb{0}\in S$ or there exists $v\in S$ such that $v\in\span(S\setminus\{v\})$.
\bpr
If $\mb{0}\in S$, its coefficient can be any. Otherwise, pick a nonzero $c_k$, the corresponding vector can be written as a linear combination of others.
\epr
\subsubsection{Linear independent set as minimum span generating set by order of inclusion}
A subset $X$ of a vector space is linearly dependent iff there exists a proper subset of $X$ whose span is the span of $X$.
\bpr
If: Let $X=\{x_1,\ldots,x_n\}$. There exist not all zero $c_i$ such that
\[\sum_{i=1}^nc_ix_i=0\]
Take any $k$ such that $c_k\ne 0$, we can write $x_k$ in terms of other $x_i$, thus $x_k$ is in the span of $X\setminus\{x_k\}$, thus $X\setminus\{x_k\}$ is a proper subset of $X$ whose span is the span of $X$.

Only if: Suppose $Y$ is proper subset of $X$ such that the span of $Y$ is the span of $X$. Choose $x\in X\setminus Y$. By the definition, there exists $y_1,\ldots,y_m\in Y$ such that $x$ is a linear combination of them.
\epr
\subsubsection{Spanning set element exchange theorem}
Let $S$ be a set of vectors and $v$ be a nonzero vector.
\[v\in\span S\implies\exists w\in S\colon w\in\span((S\setminus\{w\}\cup\{v\}))\]
\bpr
Write $v$ as a linear combination of elements of $S$. Choose one of them with nonzero coefficient, $w$, which must exist since $v$ is nonzero, $w$ can be written as a linear combination of elements in $S\setminus\{w\}\cup\{v\}$.
\epr
Let $S$ be a set of linearly independent vectors and $v$ be a nonzero vector.
\[v\in\span S\iff\exists w\in S\colon w\in\span((S\setminus\{w\}\cup\{v\}))\]
\bpr
Write $w$ as a linear combination of elements of elements in $(S\setminus\{w\}\cup\{v\}$. The $v$ term must have nonzero coefficient since otherwise $w\in\span(S\setminus\{w\})$ and $S$ is not linearly independent. Rearrange the linear combination gives $v$ as a linear combination of elements in $S$.
\epr
\subsubsection{Steinitz exchange lemma}
Let $U$ be a set of $m\in\bbN$ linearly independent vectors and $W$ be a set of $n\in\bbN$ vectors such that $U\subseteq\span W$. Then, $m\le n$, and, there exists a subset $W'$ of $W$ such that $\abs{W'}\le n-m$ and $\span(U\cup W')=\span W$.
\bpr
First statement: Write both $U$ and $W$ as column vectors of arbitrary order. By $U\subseteq\span W$, there exists $A\in F^{m\times n}$ such that $U=AW$. Because $U$ is linearly independent, $A$ must have row rank $m$. Since $\rank(A)\le\min(m,n)$, $m\le n$.

Second statement: Let $W_0=W$. For each $u_i\in U$, write $u_i$ as a linear combination of elements of $W_{i-1}$. Choose one of them with nonzero coefficient, $w_i$, which must exist since $U$ is linearly independent, $w_i$ can be written as a linear combination of elements in $W_i\coloneq W_{i-1}\setminus\{w_i\}\cup\{u_i\}$.
\epr
\subsubsection{Basis (基)}
A basis of a vector space $V$ over a field $F$ is a set of linearly independent vectors $\{v_1,v_2,\ldots\}\subseteq V$ such that for any vector $v\in V$, there exist unique coefficients $c_1,c_2,\ldots\in F$ such that:
\[\sum_{i=1}c_iv_i=v.\]
Basis for a vector space containing only zero vector is undefined.
\subsubsection{Orientation}
The orientation of a vector space is the arbitrary choice of which ordered bases are "positively" oriented and which are "negatively" oriented.
\subsubsection{Dimension (維度)}
The dimension of a vector space $V$, denoted as $\dim V$, is defined as the cardinality of a basis of it, which is unique for finite-dimensional vector space (proven below), and is unique for infinite dimension if axiom of choice is assumed (not proven here).
\bpr
Assume that $(u_1,\ldots,u_m)$ and $(v_1,\ldots,v_n)$ are two bases of vector space $V$. Since the existence of the coefficients and the fact that $u_1,\ldots,u_m,v_1,\ldots,v_n$ are nonzero, $u_1,\ldots,u_m\in\span\{v_1,\ldots,v_n\}$ and $v_1,\ldots,v_n\in\span\{u_1,\ldots,u_m\}$. By Steinitz exchange lemma, $m\le n$ and $n\le m$. Thus $m=n$.
\epr
The dimension of an affine space $V$, denoted as $\dim V$, is defined as the dimension of its associated vector space.
\subsubsection{Basis extension theorem}
Let $V$ be a vector space such that there exists a basis of $V$ that contains $n\in\bbN$ nonzero vectors. For any linearly independent set of finitely many vectors $\{v_1,\ldots,v_m\}$ in $V$, if $m=n$, it is a basis; otherwise, there exists a set of $n-m$ vectors $\{v_{m+1},\ldots,v_n\}$ in $V$ such that $\{v_1,\ldots,v_n\}$ is linearly independent, and any such $v_1,\ldots,v_n$ is a basis of $V$.
\bpr
For $m=n$, we claim that $\span\{v_1,\ldots,v_n\}=V$, thus $\{v_1,\ldots,v_n\}$ is a basis. For $m<n$, since $S\coloneq\span\{v_1,\ldots,v_m\}$ is of dimension $m<n$, choose any $v_{m+1}\in V\setminus S$. Assume for contradiction that $\{v_1,\ldots,v_{m+1}\}$ is linearly dependent, we can write $v_{m+1}$ as a linear combination of $v_1,\ldots,v_m$, which means $v_{m+1}\in S$, which contradicts $v_{m+1}\in V\notin S$. Thus $\{v_1,\ldots,v_{m+1}\}$ is linearly independent. We can repeat until $n$ such that $\span\{v_1,\ldots,v_n\}=V$.
\epr
\sssc{Coordinate notation}
For an $n$-dimensional vector space $V$ over field $F$ and a chosen ordered basis $B=(v_1,\ldots,v_n)$, every vector $v\in V$ can be represented as an $n$-scalar ordered tuple $(x_1,x_2,\ldots,x_n)$ where $x_1,\ldots,x_n\in F$ are such that
\[v=\sum_{j=1}^mx_jv_j\]
We call $x_1,\ldots,x_n$ components for a vector and coordinates for the vector space relative to $B$, and denote $v$ as
\[(x_1,x_2,\ldots,x_n)\]
or
\[\langle x_1,x_2,\ldots,x_n\rangle\]
called coordinate notation in terms of $B$.
\subsubsection{Bilinear form (雙線性形式)}
Let $V$ be a vector space of dimension $n$ over a field $\mathbb{K}$. A map $B\colon V\times V\rightarrow \mathbb{K}$ is a bilinear form on the space if:
\begin{enumerate}
\item $\forall u,v,w\in V\colon B(u+v,w)=B(u,w)+B(v,w)$.
\item $\forall u,v\in V,\lambda\in \mathbb{K}\colon B(\lambda u,v)=\lambda B(u,v)$.
\item $\forall u,v,w\in V\colon B(u,v+w)=B(u,v)+B(u,w)$.
\item $\forall u,v\in V,\lambda\in \mathbb{K}\colon B(u,\lambda v)=\lambda B(u,v)$.
\end{enumerate}
\subsubsection{Symmetric bilinear form (對稱雙線性形式)}
Let $V$ be a vector space of dimension $n$ over a field $\mathbb{K}$. A map $B\colon V\times V\rightarrow \mathbb{K}$ is a symmetric bilinear form on the space if:
\begin{enumerate}
\item $\forall u,v\in V\colon B(u,v)=B(v,u)$.
\item $\forall u,v,w\in V\colon B(u+v,w)=B(u,w)+B(v,w)$.
\item $\forall u,v\in V,\lambda\in \mathbb{K}\colon B(\lambda u,v)=\lambda B(u,v)$.
\end{enumerate}
\subsubsection{Topological vector space (TVS) (拓樸向量空間)}
A topological vector space $X$ is a vector space over a topological field $\mathbb{K}$, such that vector addition $X\times X\to X$ and scalar multiplication $\mathbb{K}\times X \to X$ are continuous functions.
\subsubsection{Functional (泛函)}
A functional is a function from a vector space into the field of real or complex numbers. For a vector space $X$, $x\in X$ and a functional on $f$ on $X$, $f(x)$ is also denoted as $\langle f,x\rangle$ or $fx$.
\subsubsection{Seminorm (半範數)}
Given a vector space $X$ over a ordered field $F$, a seminorm on $X$ is a real-valued function $p\colon X\to\mathbb{R}$ with the following conditions, where $|s|$ denotes the usual absolute value of a scalar $s$:
\begin{enumerate}
\item Subadditivity/Triangle inequality:\[\forall x,y\in X\colon p(x+y)\leq p(x)+p(y).\]
\item Absolute homogeneity: \[\forall s\in\mathbb{R},x\in X\colon p(sx)=|s|p(x).\]
\end{enumerate}
These conditions implies that:
\begin{enumerate}
\item Non-negativity: $\forall x\in X\colon p(x)\geq 0$.
\item $p(0)=0$.
\end{enumerate}
\subsubsection{Norm (範數/模長)}
A norm on $X$ is a seminorm $p\colon X\to\mathbb{R}$ with the following properties:
\begin{itemize}
\item Positive definiteness/Positiveness/Point-separating: \[\forall x\in X\colon p(x)=0\implies x=0\]
\end{itemize}
A vector $x$ such that $p(x)=1$ is called an unit vector.
\subsubsection{Standard basis}
For a normed finite-dimensional vector space, a basis consisting of unit vectors must exist.
\bpr
Take a basis and divide every element with its norm.
\epr
When such a basis is chosen as the standard basis, its elements are called standard (unit/basis) vectors, and are often denoted with $\mb{e}_1,\mb{e}_2,\ldots$ if the basis is ordered.
\subsubsection{Sublinear function (亞線性函數) or quasi-seminorm (準半範數)}
Let $p\colon X\to\mathbb{R}$ be a function on a vector space $X$ over the field $\mathbb{K}$, which is either $\mathbb{R}$ or $\mathbb{C}$. $p$ is called a sublinear function or quasi-seminorm if it satisfies the following conditions:
\begin{enumerate}
\item Subadditivity/Triangle inequality:\[\forall x,y\in X\colon p(x+y)\leq p(x)+p(y).\]
\item Nonnegative homogeneity: \[\forall x\in X,r\geq 0\colon p(rx)=rp(x).\]
\end{enumerate}
\subsubsection{Operator norm or induced norm}
For normed vector spaces $V,W$, the vector space of linear operators $L(V,W)$ is equipped with a norm $\|\cdot\|_L$ called operator norm or induced norm defined as, for $A\in L(V,W)$,
\[\|A\|_L=\sup_{\|v\|_V\le 1}\|A(v)\|_W.\]
\subsubsection{Absolutely continuous function (絕對連續函數)}
Function $f\colon I\subseteq\mathbb{R}\to\mathbb{R}$ is absolutely continuous on $I$ if for every positive number $\varepsilon$, there is a positive number $\delta$ such that, for every finite sequence of pairwise disjoint sub-intervals $(x_k,y_k)$ of $I$ with $x_k<y_k\in I$ and cardinality $N$,
\[\sum_{k=1}^N(y_k-x_k)<\delta\implies\sum_{k=1}^N|f(y_k)-f(x_k)|<\varepsilon.\]
The collection of all absolutely continuous functions on $I$ is denoted $\operatorname {AC} (I)$.
\subsubsection{Bounded linear operator (有界線性運算子)}
A linear operator $T\colon X\to Y$ between two normed spaces $X$ and $Y$ is called bounded iff there exists some $M>0$ such that
\[\forall x\in X\colon\norm{T(x)}_Y\leq M\norm{x}_X.\]
The smallest such $M$ is called the operator norm of $L$ and denoted by $\norm{T}$.
\sssc{Linear operator is bounded iff Lipschitz continuous iff continuous iff continuous at 0}
A linear operator $T\colon X\to Y$ between two normed spaces $X$ and $Y$ is bounded iff Lipschitz continuous iff continuous iff continuous at 0.
\bpr
Lipschitz continuous implies continuous, and continuous implies continuous at 0, so we will show that bounded implies Lipschitz continuous, and continuous at 0 implies bounded.

Bounded implies Lipschitz continuous: Since $T$ is linear,
\[\norm{T(x)-T(y)}_Y=\norm{T(x-y)}_Y\le M\norm{x-y}_X\]
Continuous at 0 implies bounded: Since $T$ is linear, $T(0)=0$. There exists $\delta>0$ such that
\[\norm{x}_X<\delta\implies\norm{T(x)}_Y<1\]
For any $x\ne 0$, define
\[u=\frac{\delta}{2\norm{x}_X}x\]
Then
\[\norm{u}_X=\frac{\delta}{2}<\delta\]
Since $T$ is linear,
\[T(u)=\frac{\delta}{2\norm{x}_X}T(x)\]
Hence,
\[\frac{\delta}{2\norm{x}_X}\norm{T(x)}_Y=\norm{T(u)}_Y<1\]
which implies
\[\norm{T(x)}_Y<\frac{2}{\delta}\norm{x}_X\]
The same inequality also holds for $x=0$. Hence $T$ is bounded with $M=\frac{2}{\delta}$.
\epr
\sssc{Function space as vector space}
Let $S$ be a set and $V$ be a vector space over some field $F$. Then, the set of all functions with domain $S$ and codomain $V$ is a vector space over $F$.
\subsubsection{Uniform norm, sup norm, supremum norm, Chebyshev norm, infinity norm, or max norm}
Uniform norms are defined on the function space of bounded functions valued in a normed space. The uniform norm of bounded function $f\colon X\to(Y,\|\|_Y)$ is defined by
\[\|f\|_{\infty}=\sup_{x\in X}\|f(x)\|_Y.\]
\subsubsection{Hahn–Banach theorem (哈恩-巴拿赫定理)}
Let $p\colon X\to\mathbb{R}$ be a sublinear functional on a vector space $X$ over the field $\mathbb{K}$, which is either $\mathbb{R}$ or $\mathbb{C}$. If $f\colon M\to\mathbb{K}$ is a linear functional on a vector subspace $M$ such that
\[\forall m\in M\colon f(m)\leq p(m),\]
then there exists a linear functional $F\colon X\to\mathbb{K}$ such that
\[\forall m\in M\colon F(m)=f(m),\]
\[\forall x\in X\colon F(x)\leq p(x).\]
\subsubsection{p-norm (p-範數)}
The $p$-norm of a $n$-dimensional vector $x=((x_i)_{i=1}^n)$ is defined to be
\[\|x\|_p=\qty(\sum_{i=1}^n|x_i|^p)^{\frac{1}{p}},\quad p\in\mathbb{N}_{\infty},\]
where the $\infty$-norm is
\[\|x\|_{\infty}=\max_{i=1}^n|x_i|.\]
The 1-norm is also called Manhattan norm (曼哈頓範數); the 2-norm is also called Euclidean norm (歐幾里德範數).
\subsubsection{Locally convex topological vector space (LCTVS) or locally convex space}
A locally convex topological vector space (LCTVS) $(X,\mathcal{T})$ is a TVS whose topology is generated by a family of seminorms $P$ on it such that:
\[\bigcap_{p\in P}\{x\colon p(x)=0\}=\{0\}\]
That is, $\mathcal{T}$ is generated by a basis $\{U_{\epsilon,x}\}_{\epsilon\in\mathbb{R}_{\geq 0},x\in X}$ of neighborhoods defined to be:
\[U_{\epsilon,0}=\{y\in X\colon \forall p\in P\colon p(y)<\epsilon\}.\]
\[U_{\epsilon,x}=\{y\in X\colon y-x\in U_{\epsilon,0}\}.\]
This definition implies that a LCTVS is necessarily a Hausdorff space (T2 space).
\subsubsection{Banach space (巴拿赫空間)}
A Banach space is a normed vector space $(X,\|{\cdot }\|)$ with induced metric $d$
\[d(x,y)=\|x-y\|,\quad \forall x,y\in X\]
that is a Complete metric space (aka Cauchy space).
\subsubsection{Dot product (點積)}
The dot product of two vectors $\mathbf{a}=(a_1,a_2,\dots,a_n)$ and $\mathbf{b}=(b_1,b_2,\dots,b_n)$ is defined as:
\[\mathbf{a}\cdot\mathbf{b}=\sum_{i=1}^n a_ib_i.\]
Two nonzero vector are called orthogonal or perpendicular iff their dot product is zero.
\subsubsection{Inner product space (內積空間)}
A bar over a scalar denotes the complex conjugate of it. An inner product space is a vector space $V$ over the real or complex field $F$ together with an inner product, that is, a map
\[\langle \cdot ,\cdot \rangle \colon V\times V\to F\]
that satisfies the following three properties for all vectors $x,y,z\in V$ and all scalars $a,b\in F$.
\bit
\item Conjugate symmetry: 
\[\langle x,y\rangle =\ol{\langle y,x\rangle }.\]
This implies that $\langle x,x\rangle$ is a real number.
\item Linearity in the first argument:
\[ \langle ax+by,z\rangle =a\langle x,z\rangle +b\langle y,z\rangle .\]
\item Positive-definiteness: if $x$ is not zero, then $\langle x,x\rangle >0$.
\eit
Every inner product space has an induced norm called canonical norm and is often assumed to be equipped defined by
\[\norm{x}=\sqrt{\langle x,x\rangle}\]
The inner product space over $\bbC$ is also called unitary space.
\subsubsection{Inner product is bilinear form}
An inner product is a bililinear form. TODO
\bpr
TODO
\epr
\subsubsection{Orthonormal basis (正交基)}
An orthonormal basis for an inner product space $V$ with finite dimension is a basis for $V$ such that for $e,f$ in it:
\[\langle e,f \rangle = \begin{cases}1, & \text{if } e=f,\\0, & \text{if } e\neq f.\end{cases}\]
\subsubsection{Balanced set (平衡集), circled set, or disk set}
Let $X$ be a vector space over a field $\mathbb{K}$ with an absolute value function $|\cdot |$. A subset $S$ of $X$ is called a balanced set or balanced if:
\[\forall a\in\mathbb{K}\text{ s.t. }|a|\leq 1\colon aS\subseteq S.\]
\subsubsection{Well-defined (定義良好)/Unambiguous (不模糊的)}
A well-defined expression or unambiguous expression is an expression whose definition assigns it a unique interpretation or value. Otherwise, the expression is said to be not well defined, ill defined or ambiguous.
\subsubsection{Affine map (仿射映射)}
Given two affine spaces $A$ and $B$ whose associated vector spaces are $\mathscr{A}$ and $\mathscr{B}$, an affine map from $A$ to $B$ is a map $f\colon A\to B$ such that
\[\begin{aligned}
&\mathscr{f}\colon\mathscr{A}\to\mathscr{B}\\
&b-a\mapsto f(b)-f(a)
\end{aligned}\]
is a well-defined linear map.
\subsubsection{Affine transformation (仿射變換)}
An affine transformation is a bijection from an affine space onto itself that is an affine map.
\subsubsection{Von Neumann bounded or bounded}
Suppose $X$ is a topological vector space over a field $\mathbb{K}$. A subset $B$ of $X$ is called von Neumann bounded (or bounded) in $X$ if for every neighborhood $V$ of the origin, there exists a real $r>0$ such that $B\subseteq sV$ for all scalars $s$ satisfying $|s|\geq r$.
\subsubsection{Convex set (凸集)}
Let $S$ be a vector space or an affine space over some ordered field. A subset $C$ of $S$ is convex if, for all $x$ and $y$ in $C$, the line segment connecting $x$ and $y$ is included in $C$.
\subsubsection{(Algebraic) dual (對偶) (vector) (space)}
Given any vector space $V$ over a field $F$, the dual space $V^*$ is defined as the set of all linear maps $\varphi\colon V\to F$. The dual space $V^*$ itself thus becomes a vector space over $F$ when equipped with an addition and scalar multiplication satisfying that for any $\varphi ,\psi \in V^*$, for any $x\in V$, for any $a\in F$,
\[(\varphi+\psi)(x)=\varphi (x)+\psi (x),\]
\[(a\varphi )(x)=a\left(\varphi (x)\right).\]
Element in dual is called covector.
\subsubsection{Continuous dual (space)}
Given any vector space $V$ over a field $F$, the continuous dual space $V'$ is defined as the set of all continuous linear maps $\varphi\colon V\to F$. The continuous dual space $V'$ itself thus becomes a vector space over $F$ when equipped with an addition and scalar multiplication satisfying that for any $\varphi ,\psi \in V'$, for any $x\in V$, for any $a\in F$,
\[(\varphi+\psi)(x)=\varphi (x)+\psi (x),\]
\[(a\varphi )(x)=a\left(\varphi (x)\right).\]
\subsubsection{Transpose, dual map, or adjoint of linear map}
For two vector spaces $X,Y$ and a linear map $T\colon X\to Y$, the transpose, dual map, or adjoint of $T$ is defined as the unique map $^tT\colon Y^*\to X^*$, which is a linear map, such that for all $x\in X$, for all $y\in Y^*$,
\[\langle y,T(x)\rangle=\langle^tT(y),x\rangle.\]
If $T$ is continuous, the restriction of $^tT$ to $W'$ has range $V'$ and is continuous.
\subsubsection{Strong dual (space)}
A strong dual space is a continuous dual space of a TVS equipped with the strong (dual) topology, also called topology of uniform convergence on bounded subsets. Suppose $X$ is a TVS over field $F$ and $\mathcal{B}$ is a family of bounded subsets of $X$ such that every bounded subset of $X$ is a subset of some $B\in\mathcal{B}$. Define the family $T$ of seminorms $\|\cdot\|_B$ on $X'$ for all $x'\in X'$ as
\[T=\{\|x'\|_B\coloneq\sup_{x\in B}|x'(x)|\mid B\in\mathcal{B}\}.\]
The strong dual topology is defined as the topology generated by $T$, and thus $B'$ is a LCTVS.
\subsubsection{Orthogonal complement (正交補)}
An orthogonal complement of a subspace $W$ of a vector space $V$ equipped with a bilinear form $B$ is the set $W^{\perp }$ of all vectors in $V$ that are orthogonal to every vector in $W$, of which $v\in V$ is orthogonal to $w\in V$ if and only if $B(v,w)=0$.
\subsubsection{Affine combination}
An affine combination is a linear combination of points where all coefficients sum up to 1.
\subsubsection{Convex combination}
A convex combination is a linear combination of points where all coefficients are non-negative and sum up to 1.
\subsubsection{Coordinate system (座標/坐標系統)}
A coordinate system is a system that uses one or more numbers, or coordinates, to uniquely determine the points on a manifold.
\subsubsection{Projection (投影)}
A (vector) projection onto a subset $U$ of a vector space $V$ is a linear operator $P\colon V\to U$ such that $P^2=P$, denoted as $\proj_UP$. The property $P^2=P$ is called idempotency (冪等性). Its norm is called scalar projection denoted as $\comp_UP$.
\subsubsection{Orthographic projection (正射影)}
An orthographic on a subset $U$ of a vector space $V$ is a projection $P\colon V\to U$  such that 
\[ \forall v\in V\colon v - P(v)\in U^*\]
where $U^*$ is the orthogonal complement of \( U \).
\subsubsection{Hilbert space (希爾伯特空間)}
A Hilbert space is a real or complex inner product space that is also a complete metric space with respect to the distance function induced by the inner product.
\subsubsection{Angle (夾角) between two vectors}
Given two vectors $u,v$ in an inner product space, the angle $\theta$ between than is defined as
\[\theta=\arccos(u\cdot v).\]
\subsubsection{Principal angles (主夾角) between two subspaces}
Given two subspaces $U,W$ of an inner product space with $\dim U=k\leq\dim W=m$, there exists then a sequence of $k$ angles $0\leq\theta_1\leq\theta_2\leq\ldots\leq\theta_k\leq\frac{\pi}{2}$ called the principal angles between $U$ and $W$, the first one defined as
\[\theta_1\coloneq\min\left\{\arccos\left(\frac{|\langle u,w\rangle|}{\|u\|\|w\|}\right)\middle |u\in U\land w\in W\right\}=\angle(u_1\in U,w_1\in W),\]
The vectors $u_1$ and $w_1$ are the corresponding principal vectors.

The other principal angles and vectors are then defined recursively via
\[\theta_i\coloneq\min\left\{\arccos\left(\frac{|\langle u,w\rangle|}{\|u\|\|w\|}\right)\middle |u\in U\land w\in W\land \forall j\in\mathbb{N}\land j<i\colon\langle u_i,u_j\rangle=\langle w_i,w_j\rangle=0\right\}=\angle(u_i\in U,w_i\in W).\]
\subsubsection{Flat, hyperplane, or affine subspace (仿射子空間)}
A flat or an affine subspace of an affine space is a subspace of it that is an affine space.
\subsubsection{Hypersurface}
A hypersurface is a $(n-1)$-dimensional manifold that is a subspace of an $n$-dimensional manifold.
\subsubsection{Parallel (平行)}
Two nonzero vectors $u,v$ in a vector space over field $F$ are called to be parallel if there exists a scalar $x\in F$ such that $xu=v$, denoted as $u\parallel v$.

Two affine subspaces are called to be parallel if the direction spaces of them $D,E$ satisfies:
\[D\subseteq E\lor E\subseteq D,\]
denoted as $D\parallel E$.
\subsubsection{Coplanar}
Multiple elements of a vector or affine space are called coplanar if they are in the same 2D affine subspace.
\subsubsection{Elementwise multiplication}
For two vectors consists of scalar components $\mb{u}=(u_1,\ldots,u_n)$ and $\mb{v}=(v_1,\ldots,v_n)$, their elementwise multiplication, denoted as $\mb{u}\odot\mb{v}$, is defined as
\[\mb{u}\odot\mb{v}=(u_1v_1,\ldots,u_nv_n).\]
\subsubsection{Unit vector (單位向量) or normalized vector}
A unit vector or a normalized vector in a normed vector space is a vector of norm 1, denoted by a lowercase letter with a circumflex or hat, as in $\hat{\mathbf{v}}$ (pronounced "v-hat"). The unit vector or normalized vector $\hat{\mathbf{u}}$ of a nonzero vector $\mathbf{u}$ in a normed vector space equipped with a norm $\|\cdot\|$ is defined as:
\[\hat{\mathbf{u}}=\frac{\mathbf{u}}{\|\mathbf{u}\|}.\]
\subsubsection{Affine independence}
For an affine space $U$ associated with vector space $V$, $u_0,u_1,\ldots,u_n\in U$ are affinely independent iff $(u_1-u_0),(u_2-u_0),\ldots,(u_n-u_0)$ are linearly independent in $V$.
\subsubsection{Simplex}
A $k$-simplex or $k$-dimensional simplex in an affine space $U$ is a set
\[\{\sum_{i=0}^k\theta_iu_i\mid\sum_{i=0}^k\theta_i=1\land\theta_0,\theta_1,\ldots,\theta_k\in\bbR_{\ge 0}\},\]
where
\[u_0,u_1,\ldots,u_k\in U\]
are affine independent and called vertices of the simplex.
\subsubsection{Gram matrix, Gramian matrix, or Gramian}
For $n$ vectors $v_1,v_2,\ldots,v_n$ in an inner product space, the Gram matrix, Gramian matrix, or Gramian $G$ is defined as an $n\times n$ matrix where
\[G_{ij}=\langle v_i,v_j\rangle.\]
The Gram determinant or Gramian is the determinant of the Gram matrix.
\subsubsection{Support of functional}
For a function $f\colon U\subseteq\bbR^n\to\bbR$ (or $\bbC$), the support of it, denoted as $\supp(f)$, is defined as
\[\supp(f)=\closure(\{x\in U\mid f(x)\neq 0\}).\]
If $\supp(f)$ is a compact set, $f$ is called with compact support. If $\supp(f)$ is a compact subset of a set $V\subseteq U$, $f$ is called with compact support on $V$.
\subsubsection{Euclidean vector space}
A Euclidean vector space is a finite-dimensional inner product space over the real numbers, denoted as $\bbR^n$ for dimension $n$. The norm-induced metric $d(v,w)=\|v-w\|$ is called the Euclidean distance
\subsubsection{Euclidean affine space (歐幾里德仿射空間) or Euclidean space (歐幾里得空間)}
A Euclidean space, also known as Euclidean affine space, is an affine space over $\mathbb{R}$ such that the associated vector space is a Euclidean vector space, and is equipped with the same norm and metric as the associated Euclidean vector space.
\subsubsection{Affine hull (仿射包)}
The affine hull $\operatorname{aff} (S)$ of a set of points $S$ in a Euclidean space is the smallest affine space that contains all the points in $S$, namely, the set of all affine combinations of the points in $S$, i.e.:
\[\operatorname{aff}(S)=\left\{\sum _{i=1}^k\alpha _ix_i\colon k>0,x_i\in S,\alpha _i\in\mathbb{R},\sum _{i=1}^k\alpha _i=1\right\}\]
\subsubsection{Convex hull (凸包)}
The convex hull $\operatorname{conv} (S)$ of a set of points $S$ in a Euclidean space is the smallest convex set that contains all the points in $S$, namely, the set of all convex combinations of points in $S$, i.e.:
\[\operatorname{conv}(S)=\left\{\sum _{i=1}^k\alpha _ix_i\colon k>0,x_i\in S,\alpha _i\in [0,1],\sum _{i=1}^k\alpha _i=1\right\}\]
\subsubsection{Cartesian coordinate system (笛卡爾座標系統) or rectangular coordinate system (直角座標系統)}
The Cartesian or rectangular coordinate system is a coordinate system in a Euclidean vector space $E$ determined by an orthogonal basis of unit vectors such that the inner product of two vectors in $E$ equals to the dot product of them, and that the norm of $E$ is 2-norm.
\subsubsection{Lattice points (格子點)}
In a Cartesian coordinate system, vectors whose components are all integers are called lattice points.
\subsubsection{Curvilinear coordinate system}
A curvilinear coordinate system is a coordinate system for Euclidean space such that there exists an invertible map from a Cartesian coordinate system to it.
\subsubsection{Standard, unit, or probability simplex}
The standard, unit, or probability $k$-simplex is the simplex in $\bbR^{k+1}$ whose vertices are the unit vectors of each axis, that is, the set
\[\{(p_0,p_1,\ldots,p_k)\mid\sum_{i=0}^kp_i=1\land p_0,p_1,\ldots,p_k\in\bbR_{\ge0}\}.\]
\subsubsection{Rectangle}
An open rectangle in $\mathbb{R}^n$ is the product of $n$ open intervals.

An closed rectangle in $\mathbb{R}^n$ is the product of $n$ closed intervals.

Rectangle may refer to either, or product of $n$ intervals, or product of $n$ intervals satisfying certain constraints, depending on context, and often the case they make no non-trivial difference if not specified.
\subsection{Tensor}
\subsubsection{Tensor}
Let $V$ be a vector space and $V^*$ be its dual. A tensor of type $(r,s)$ is a multilinear map
\[T\colon (V^*)^r\times V^s\to\bbR,\]
where $r$ is the number of contravariant slots, which take covectors in $V^*$, and $s$ is the number of covariant slots, which take vectors in $V$. The order or rank of it is $r+s$.

This map
\[T\colon (V^*)^r\times V^s\to\bbR\]
is also a map
\[T\colon (V^*)^q\times V^t\to\{S\colon (S\colon(V^*)^{r-q}\times V^{s-t}\to\bbR)\}\]
for any $q\in\bbZ\cap[0,r]$ and $t\in\bbZ\cap[0,s]$.

A scalar is a tensor of type $(0,0)$. A vector is a tensor of type $(1,0)$, which maps a covector to a real number, thus also called a contravariant tensor. A covector is a tensor of type $(0,1)$, which maps a vector to a real number, thus also called a covariant tensor.. A matrix as a linear map $V\to V$ is a tensor of type $(1,1)$. A matrix as a bilinear form $V\times V\to\bbR$ is a tensor of type $(0,2)$. A matrix as $V^*\times V^*\to\bbR$ is a tensor of type $\bbR$.

It's conventional that a vector is a column vector and a covector is a row vector. Tensor convention follows.

A tensor can be represented as an array whose dimension is its rank when a basis is chosen.
\subsubsection{Upper and lower indices}
Dummy indices such as $\mu$ and $\nu$ are used to indicate contravariant and covariant vectors and tensors.

We use a superscript of $r$ dummy indices and a subscript of $s$ dummy indices to indicate a tensor of type $(r,s)$, i.e., $(V^*)^r\times V^s\to\bbR$ for vector space $V$, such as $T^{\mu_1\ldots\mu_r}_{\nu_1\ldots\nu_s}$. Each $i$th component of each upper index $\mu$ is denoted by replacing $\mu$ with $i$. Each $i$th component of each lower index $\mu$ is denoted by replacing $\mu$ with $i$. For vectors, that is, $A^{\mu}=(A^0,A^1,A^2,A^3)$. For covectors, that is, $A_{\mu}=(A_0,A_1,A_2,A_3)$.

Each upper index $\mu$ is fixed to a summation $\sum_{\mu=0}^{\dim V}$ when it takes an input of the same upper index $\mu$, i.e.
\[T^{\mu\ldots_1}_{\ldots_2}U^{\mu\ldots_3}_{\ldots_4}=\sum_{\mu=0}^{\dim V}T^{\mu\ldots_1}_{\ldots_2}U^{\mu\ldots_3}_{\ldots_4}.\]
Each lower index $\mu$ is fixed to a summation $\sum_{\mu=0}^{\dim V}$ when it takes an input of the same lower index $\mu$, i.e.
\[T^{\ldots_1}_{\mu\ldots_2}U^{\ldots_3}_{\mu\ldots_4}=\sum_{\mu=0}^{\dim V}T^{\ldots_1}_{\mu\ldots_2}U^{\ldots_3}_{\mu\ldots_4}.\]
The indexing replaces the order of them in a formula and makes the multiplication notation of tensors work without extra notations.
\subsection{Two dimensional Euclidean vector space}
\subsubsection{Quadrant (象限)}
\RNum{1}(+, +), \RNum{2}(-, +), \RNum{3}(-, -), \RNum{4}(+, -).
\subsection{Three-dimensional Euclidean vector space}
\subsubsection{Right-hand rule (右手定則)}
Right-hand rule is a common used orientation for a three-dimensional space. In a right-handed Cartesian coordinate system, if you point the thumb of your right hand in the positive $x$-axis direction and your index finger in the positive $y$-axis direction, then your middle finger (extended perpendicularly from the palm) will point in the positive $z$-axis direction. If not specified otherwise, right-handed Cartesian coordinate system is usually used.
\subsubsection{Octant (卦限)}
\RNum{1}(+, +, +), \RNum{2}(-, +, +), \RNum{3}(-, -, +), \RNum{4}(+, -, +), \RNum{5}(+, +, -), \RNum{6}(-, +, -), \RNum{7}(-, -, -), \RNum{8}(+, -, -).
\subsubsection{Cross product (叉積), external product (外積), or vector product (向量積)}
The cross product of vector $\mathbf{a}=(a_1,a_2,a_3)$ and $\mathbf{b}=(b_1,b_2,b_3)$ is defined as:
\[\mathbf{a}\times\mathbf{b}=\det\begin{pmatrix}\mathbf{i} & \mathbf{j} & \mathbf{k}\\a_1 & a_2 & a_3\\b_1 & b_2 & b_3\end{pmatrix},\]
where $\mb{i}=(1,0,0)$, $\mb{j}=(0,1,0)$, $\mb{k}=(0,0,1)$.
\subsubsection{Basic properties of cross product}
\bit
\item Self cross product is the zero vector:
\[\mb{a}\times\mb{a}=\mb{0}.\]
\item Anticommutative law (反交換律):
\[\mathbf{a}\times\mathbf{b}=-\mathbf{b}\times\mathbf{a}.\]
\item Left and right distributive over addition and subtraction:
\[\mb{a}\times(\mb{b}\pm\mb{c})=\mb{a}\times\mb{b}\pm\mb{a}\times\mb{c}.\]
\[(\mb{b}\pm\mb{c})\times\mb{a}=\mb{b}\times\mb{a}\pm\mb{c}\times\mb{a}.\]
\item Commutative and associative over scalar multiplication:
\[(r\mb{a})\times\mb{b}=a\times(r\mb{b})=r(\mb{a}\times\mb{b}).\]
\item Two vector in $\bbR^3$ are parallel iff their cross product is the zero vector.
\eit
\subsubsection{Triple product identity (三重積恆等式)/Jacobi identity (雅可比恆等式)}
\[\mathbf{A}\times(\mathbf{B}\times\mathbf{C})+\mathbf{B}\times(\mathbf{C}\times\mathbf{A})+\mathbf{C}\times(\mathbf{A}\times\mathbf{B})=0.\]
\[\mathbf{A}\cdot(\mathbf{B}\times\mathbf{C})=\mathbf{B}\cdot(\mathbf{C}\times\mathbf{A})=\mathbf{C}\cdot(\mathbf{A}\times\mathbf{B})=\begin{vmatrix}\mathbf{A}\\\mathbf{B}\\\mathbf{C}\end{vmatrix}\]
\[(\mathbf{A}\times\mathbf{B})\times\mathbf{C}=(\mathbf{A}\cdot\mathbf{C})\mathbf{B}-(\mathbf{B}\cdot\mathbf{C})\mathbf{A}.\]
\subsubsection{Scalar triple product}
\[[\mb{a},\mb{b},\mb{c}]=\det(\mb{a},\mb{b},\mb{c})=\mb{a}\cdot(\mb{b}\times\mb{c})\]
\sssc{Lami's theorem (拉密定理)}
Let $u,v,w\in\bbR^3$ be three nonzero vectors such that $u+v+w=0$. Then they are coplanar, and let $\alpha,\beta,\gamma$ be the angle directly opposite to the vectors, they can form a triangle by connecting head to end, and
\[\frac{\sin\alpha}{|u|}=\frac{\sin\beta}{|v|}=\frac{\sin\gamma}{|w|}\]
\section{Matrix (矩陣)}
\subsection{Matrix}
A matrix (plural: matrices) is a rectangular array of numbers (or other mathematical objects), called the entries (元或元素) of the matrix. Most commonly, a matrix over a field $F$ is a rectangular array of elements of $F$.

The size of a matrix is defined by the number of rows (zh-hant: 列/zh-hans: 行) and columns (zh-hant: 行/zh-hans: 列) it contains. A matrix with $m$ rows and $n$ columns is called an $m\times n$ matrix, or $m$-by-$n$ matrix, where $m$ and $n$ are called its dimensions (階). A matrix with the same number of rows and columns $n$ is called a $n$ square matrix (方陣). A matrix with an infinite number of rows or/and columns is called an infinite matrix. A matrix with exactly one row is called a row vector. A matrix with exactly one column is called a column vector.

The entry in the $i$-th row and $j$-th column, called the $i,j$th entry (第$i,j$元) or $ij$th entry (第$ij$元), is commonly denoted as $a_{i,j}$ or $a_{ij}$.

An $m\times n$ matrix $\mathbf{A}$ is commonly denoted as:
\[\mathbf{A}=\begin{bmatrix}a_{11} & a_{12} & \ldots & a_{1n}\\a_{21} & a_{22} & \ldots & a_{2n}\\\vdots & \vdots & \ddots & \vdots\\a_{m1} & a_{m2} & \ldots & a_{mn}\end{bmatrix},\]
\[\mathbf{A}=\begin{pmatrix}a_{11} & a_{12} & \ldots & a_{1n}\\a_{21} & a_{22} & \ldots & a_{2n}\\\vdots & \vdots & \ddots & \vdots\\a_{m1} & a_{m2} & \ldots & a_{mn}\end{pmatrix},\]
\[\mathbf{A}=[a_{ij}]_{1\leq i\leq m,1\leq j\leq n},\]
\[\mathbf{A}=(a_{ij})_{1\leq i\leq m,1\leq j\leq n},\]
\[\mathbf{A}=[a_{ij}]_{m\times n},\]
or
\[\mathbf{A}=(a_{ij})_{m\times n}.\]

The $i,i$th entries are called diagonal entries (對角線元), and other entries are called off-diagonal or non-diagonal entry. The diagonal line of diagonal entries is called main diagonal.

A matrix with only one row is called a row matrix or row vector, and a matrix with only one column is called a column matrix or column vector.

The $i$th row of the matrix $\mathbf{A}$ is commonly denoted as $\mathbf{A}_{i,:}$ or $\mb{A}_i$, and the $j$th column is commonly denoted as $\mathbf{A}_{:,j}$.

The set of all $m\times n$ matrices over $\mathbb{F}$ is denoted as $\mathbb{F}^{m\times n}$.
\subsection{Basic operations}
\subsubsection{Submatrix}
A submatrix of a matrix $M$ is obtained by deleting from $M$ entire rows, entire columns, or both.
\subsubsection{Matrix equality}
We say that two matrices $\mathbf{A}=[a_{ij}]_{m\times n}$ and $\mathbf{B}=[b_{ij}]_{m\times n}$ are equal if
\[\forall 1\leq i\leq m,1\leq j\leq n\colon a_{ij}=b_{ij}.\]
\subsubsection{Matrix addition}
The sum of two $m\times n$ matrices $\mathbf{A}$ and $\mathbf{B}$ is calculated entrywise.
\subsubsection{Scalar multiplication}
The product $c\mathbf{A}$ of a scalar $c$ is calculated by multiplying every entry of $\mathbf{A}$ by $c$.
\subsubsection{Matrix multiplication (矩陣乘法)}
Multiplication of two matrices is defined only if the number of columns of the left matrix is the same as the number of rows of the right matrix. The matrix product $\mathbf{A}\mathbf{B}$ of an $m\times n$ matrix $\mathbf{A}$ and $n\times p$ matrix $\mathbf{B}$ is the $m\times p$ matrix whose $ij$th entry is the dot product of the $i$th row of $\mathbf{A}$ and the $j$th column of $\mathbf{B}$.
\subsubsection{Outer product}
For two column vectors $u,v$, the outer product $uv^\top$, sometimes denoted as $u\otimes v$, is the matrix product $uv^\top$.
\subsubsection{Positive integer exponent}
Given a square matrix $\mathbf{A}$ and a positive integer $n$, $\mathbf{A}^n$ is defined as $\prod_{i=1}^n\mathbf{A}$.
\subsubsection{Zero matrix (零矩陣)}
A zero matrix over a field $F$ is a matrix of which all entries are the additive identity $0_F$ of $F$. Zero matrix is often denoted as $\mb{O}$ or $O$. $n\times n$ zero matrix is often denoted as $\mb{O}_n$ or $O_n$.
\subsubsection{Identity matrix (單位矩陣)}
An identity matrix over a field $F$ is a square matrix of which all diagonal entries are the multiplicative identity $1_F$ of $F$ and all off-diagonal entries are the multiplicative absorbing element $0_F$ of $F$. Identity matrix is often denoted as $\mb{I}$ or $I$. $n\times n$ identity matrix is often denoted as $\mb{I}_n$ or $I_n$.
\subsubsection{Hadamard product (哈達瑪乘積), Schur product (舒爾乘積), elementwise product, or entrywise product (逐項乘積)}
For two $m\times n$ matrices $\mb{A}$ and $\mb{B}$, the Hadamard product $ \mb{A}\odot \mb{B}$ or $\mb{A}\circ \mb{B}$ is an $m\times n$ matrix defined as:
\[(\mb{A}\odot \mb{B})_{ij}=(\mb{A})_{ij}(\mb{B})_{ij}.\]
\subsubsection{Kronecker product (克羅內克積)}
If $\mathbf{A}$ is an $m\times n$ matrix and $\mathbf{B}$ is a $p\times q$ matrix, then the Kronecker product $\mathbf{A}\otimes\mathbf{B}$ is the $pm\times qn$ block matrix:
\[(\mathbf{A}\otimes\mathbf{B})_{i,j}=a_{\ceil{\frac{i}{p}},\ceil{\frac{j}{q}}}b_{((i-1)\mod p)+1,((j-1)\mod q)+1}\]
\subsubsection{Merging matrices}
The matrix formed by merging matrices with the same number of rows left and right is denoted by writing them in columns of a row vector.

The matrix formed by merging matrices with the same number of columns up and down is denoted by writing them in rows of a column vector.
\subsubsection{Property}
For any scalars $r$ and $s$, $m\times n$ matrices $A$, $B$, $C$, $D$, zero matrix $O$, and identity matrix $I$, the following statements are true,
\[A+B=B+A\]
\[A+(B+C)=(A+B)+C\]
\[r(A+B)=rA+rB\]
\[rA+sA=(r+s)A\]
\[(rs)A=r(sA)=s(rA)\]
\[(AB)C=A(BC)\]
\[r(AB)=(rA)B=A(rB)\]
\[A(B+C)=AB+AC\]
\[(A+B)C=AC+BC\]
\[(A+B)(C+D)=AC+AD+BC+BD\]
\[A\odot B=B\odot A\]
\[A\odot(B\odot C)=(A\odot B)\odot C\]
\[A\odot(B+C)=A\odot B+A\odot C\]
\[IA=A\]
\[AI=A\]
\[A+O=A\]
\[AO=O\]
\[OA=O\]
the following statements are not necessarily true,
\[AB=BA\]
\bpr
Counterexample:
\[A=\begin{pmatrix}1&0\\0&0\end{pmatrix}\]
\[B=\begin{pmatrix}0&1\\1&0\end{pmatrix}\]
\[AB=\begin{pmatrix}0&1\\0&0\end{pmatrix}\]
\[BA=\begin{pmatrix}0&0\\1&0\end{pmatrix}\]
\epr
\[(AB=AC\land A\neq O)\Rightarrow B=C\]
\bpr
Counterexample:
\[A=\begin{pmatrix}1&0\\0&0\end{pmatrix}\]
\[B=\begin{pmatrix}1&0\\0&1\end{pmatrix}\]
\[C=\begin{pmatrix}1&0\\0&2\end{pmatrix}\]
\[AB=AC=A\]
\epr
\[(ACB=ADB\land A\neq O\land B\neq O)\Rightarrow C=D\]
\bpr
Counterexample:
\[A=\begin{pmatrix}1&0\\0&0\end{pmatrix}\]
\[B=C=\begin{pmatrix}1&0\\0&1\end{pmatrix}\]
\[D=\begin{pmatrix}1&0\\0&2\end{pmatrix}\]
\[ACB=AC=ADB=AD=A\]
\epr
\[A^2=O\Rightarrow A=O\]
\bpr
Counterexample:
\[A=\begin{pmatrix}0&1\\0&0\end{pmatrix}\]
\epr
\[AB=O\Rightarrow(A=O\lor B=O)\]
\bpr
Counterexample:
\[A=\begin{pmatrix}1&0\\0&0\end{pmatrix}\]
\[B=\begin{pmatrix}0&0\\0&1\end{pmatrix}\]
\epr
\[AB=B\Rightarrow A=I\]
\bpr
Counterexample:
\[A=B=AB=\begin{pmatrix}1&0\\0&0\end{pmatrix}\]
\epr
\subsubsection{Zero exponent}
Given an $n\times n$ matrix $\mathbf{A}$ and $n\times n$ identity matrix $\mathbf{I}$, $\mathbf{A}^0$ is defined as $\mathbf{I}$.
\subsubsection{Trace (跡)}
The trace of a matrix $\mb{A}$ is the sum of all diagonal entries of $\mb{A}$, denoted as $\tr(\mb{A})$.
\subsection{Elementary row operations}
\subsubsection{Elementary row operations (運算)}
Elementary row operations include:
\bit
\item Row addition (相加): Adding a row $j$ multiplied by a nonzero constant $c$ to another row $i\ne j$, denoted as $R_i\leftarrow R_i+cR_j$.
\item Row multiplication/scaling (加倍): Multiplying a row $i$ by a nonzero constant $c$, denoted as $R_i\leftarrow cR_i$.
\item Row interchange/swap/switch (交換): Interchanging two rows $i,j$, denoted as $R_i\leftrightarrow R_j$.
\eit
If one matrix can be transformed to another through a sequence of elementary row operations, the two matrices are called row-equivalent.
\sssc{Elementary column operations}
Elementary column operations are defined the same as row replaced with column and $R$ replaced with $C$.

If one matrix can be transformed to another through a sequence of elementary column operations, the two matrices are called column-equivalent.
\subsubsection{Elementary ow operation as matrix multiplication and elementary matrix}
Every elementary row operation on a $m\times n$ matrix $A$ can be performed by left-multiplying an appropriate $m\times m$ matrix $E$:
\bit
\item Row addition $R_i\leftarrow R_i+cR_j$: $E$ is $I_m$ with the $j,i$ entry replaced with $c$. 
\item Row multiplication $R_i\leftarrow cR_i$: $E$ is $I_m$ with the $i,i$ entry replaced with $c$. 
\item Row interchange $R_i\leftrightarrow R_j$: $E$ is $I_m$ with the row $i$ and the row $j$ interchanged.
\eit
A square matrix is called elementary matrix if it can be obtained from the identity matrix of its size with a single elementary row operations.

Let $A$ be a $m\times n$ matrix and $E$ be an $m\times m$ elementary matrix that can be obtained by performing an elementary row operation on $I_m$. Then, $EA$ is the matrix obtained by performing the same operation on $A$.
\subsubsection{Zero row, nonzero row, and leading entry}
A row whose all entries are zero is called a zero row, otherwise it is called a nonzero row. The left-most nonzero entry of a nonzero row is called leading entry.
\subsubsection{Transpose (轉置)}
The transpose of an $m\times n$ matrix $\mathbf{A}$ is the $n\times m$ matrix $\mathbf{A}^\top$:
\[(\mathbf{A}^\top)_{i,j}=a_{j,i}\]
\subsubsection{Symmetric matrix (對稱矩陣)}
$\mathbf{A}$ is called a symmetric matrix iff $\mathbf{A}=\mathbf{A}^\top$.
\subsubsection{Skew-symmetric matrix/antisymmetric matrix (反對稱矩陣)}
$\mathbf{A}$ is called a skew-symmetric matrix or an antisymmetric matrix if $\mathbf{A}=-\mathbf{A}^\top$.
\subsubsection{Transpose identities}
For an index set of matrices $\mathbf{A}\colon[1,n]\cap\bbN\to F^{p\times q};\;i\mapsto\mathbf{A}(i)$ for some fixed $p,q$,
\[\qty(\sum_{i=1}^n\mathbf{A}(i))^\top=\sum_{i=1}^n\qty(\mathbf{A}(i))^\top\]
\bpr
\[\qty(\sum_{i=1}^n\mathbf{A}(i))_{jk}=\sum_{i=1}^n\mathbf{A}(i)_{jk}\]
\[\ba
\qty(\qty(\sum_{i=1}^n\mathbf{A}(i))^{\top})_{jk}&=\sum_{i=1}^n\mathbf{A}(i)_{kj}\\
                                                 &=\sum_{i=1}^n\qty(\qty(\mathbf{A}(i))^{\top})_{jk}\\
                                                 &=\qty(\sum_{i=1}^n\qty(\mathbf{A}(i))^{\top})_{jk}
\ea\]
\epr
For two matrices $X\in F^{p\times q},Y\in F^{q\times r}$,
\[(XY)^{\top}=Y^{\top}X^{\top}\]
\bpr
\[(XY)_{jk}=\sum_{i=1}^qX_{ji}Y_{ik}\]
\[\ba
((XY)^{\top})_{jk}&=\sum_{i=1}^qX_{ki}Y_{ij}\\
                  &=\sum_{i=1}^q(X^{\top})_{ik}(Y^{\top})_{ji}\\
                  &=(Y^{\top}X^{\top})_{jk}
\ea\]
\epr
For an index set of matrices $\mathbf{A}\colon[1,n]\cap\bbN\to F^{\bbN\times\bbN};\;i\mapsto\mathbf{A}(i)$ such that $\mathbf{A}(i)\mathbf{A}(i+1)$ is defined for all $i\in[1,n)\cap\bbN$,
\[\qty(\prod_{i=1}^n\mathbf{A}(i))^\top=\prod_{i=1}^n\qty(\mathbf{A}(n-i+1))^\top\]
\bpr
Prove by mathematical induction on $n$.

Base case $n=1$:
\[\qty(\prod_{i=1}^1\mathbf{A}(i))^\top=\prod_{i=1}^1\qty(\mathbf{A}(1-i+1))^\top,\]
which is true.

Assume the statement holds for $n=k$:
\[\qty(\prod_{i=1}^k\mathbf{A}(i))^\top=\prod_{i=1}^k\qty(\mathbf{A}(k-i+1))^\top\]
For $n=k+1$, we have
\[\ba
\qty(\prod_{i=1}^{k+1}\mathbf{A}(i))^\top&=\qty(\prod_{i=1}^k\mathbf{A}(i)A(k+1))^\top\\
                                         &=\qty(A(k+1))^{\top}\qty(\prod_{i=1}^k\mathbf{A}(i))^\top\\
                                         &=\qty(A(k+1))^{\top}\prod_{i=1}^k\qty(\mathbf{A}(k-i+1))^\top\\
                                         &=\prod_{i=1}^{k+1}\qty(\mathbf{A}(k-i+1))^\top
\ea\]
\epr
\subsubsection{Symmetric and skew-symmetric decomposition}
Every real or complex square matrix can be uniquely written as a sum of a symmetric matrix $\frac{A+A^\top}{2}$ and a skew-symmetric matrix $\frac{A-A^\top}{2}$.
\bpr
If $A=S_1+K_1=S_2+K_2$ with $S_1,S_2$ symmetric and $K_1,K_2$ skew-symmetric. $S_1-S_2=K_2-K_1$ is both symmetric and skew-symmetric, which can only be the zero matrix.
\epr
\subsection{Complex matrix}
\subsubsection{(Complex）conjugate}
For a complex matrix $A$, its complex conjugate $\bar{A}$ is defined by taking complex conjugate for all its entries.
\subsubsection{Conjugate transpose (共軛轉置)/Hermitian conjugate (埃爾米特共軛)/Hermitian transpose (埃爾米特轉置)}
For a complex matrix $A$, its conjugate transpose $A^*$ is defined as
\[A^*=\bar{A}^\top\]
\subsubsection{Conjugate transpose identities}
For an index set of matrices $\mathbf{A}\colon[1,n]\cap\bbN\to\bbC^{p\times q};\;i\mapsto\mathbf{A}(i)$ for some fixed $p,q$,
\[\qty(\sum_{i=1}^n\mathbf{A}(i))^*=\sum_{i=1}^n\qty(\mathbf{A}(i))^*\]
For two matrices $X\in\bbC^{p\times q},Y\in\bbC^{q\times r}$,
\[(XY)^{*}=Y^{*}X^{*}\]
For an index set of matrices $\mathbf{A}\colon[1,n]\cap\bbN\to\bbC^{\bbN\times\bbN};\;i\mapsto\mathbf{A}(i)$ such that $\mathbf{A}(i)\mathbf{A}(i+1)$ is defined for all $i\in[1,n)\cap\bbN$,
\[\qty(\prod_{i=1}^n\mathbf{A}(i))^*=\prod_{i=1}^n\qty(\mathbf{A}(n-i+1))^*\]
\ssc{Inverse}
\subsubsection{Inverse (反方陣、逆方陣、反矩陣或逆矩陣), multiplicative inverse (乘法反方陣), or reciprocal}
The inverse of an $n\times n$ matrix $\mathbf{A}$ is an $n\times n$ matrix $\mathbf{A}^{-1}$ such that:
\[\mathbf{A}\mathbf{A}^{-1}=\mathbf{A}^{-1}\mathbf{A}=\mathbf{I}\]
If such $\mb{A}^{-1}$ exists, $\mb{A}$ is called invertible.
\subsubsection{Negative integer exponent}
Given an invertible matrix $\mathbf{A}$ and a negative integer $n$, $\mathbf{A}^n$ is defined as $(\mathbf{A}^{-1})^n$.
\subsubsection{Invertible matrix (可逆矩陣), orthogonal matrix (正交矩陣), orthonormal matrix (正則矩陣), non-singular matrix (非奇異矩陣), or non-degenerate matrix (非退化矩陣)}
An invertible matrix, orthogonal matrix, orthonormal matrix, non-singular matrix, or non-degenerate matrix is a matrix such that its inverse exists.
\subsubsection{Non-invertible matrix (不可逆矩陣), non-orthogonal matrix (非正交矩陣), non-orthonormal matrix (非正則矩陣), singular matrix (奇異矩陣), or degenerate matrix (退化矩陣)}
An non-invertible matrix, non-orthogonal matrix, non-orthonormal matrix, singular matrix, or degenerate matrix is a matrix that is not an invertible matrix.
\subsubsection{Inverse of a two-by-two matrix}
For an invertible $2\times 2$ matrix $\mathbf{A}$:
\[\mathbf{A}=\begin{bmatrix}a & b \\c & d\end{bmatrix}\implies\mathbf{A}^{-1}=
\frac{1}{\det\qty(\mathbf{A})}\begin{bmatrix}d & -b \\-c & a\end{bmatrix}\]
\subsubsection{Unitary matrix (么正矩陣或酉矩陣)}
An unitary matrix is a square matrix whose inverse equals to its conjugate transpose.
\subsubsection{Inverse identities}
A matrix is invertible iff its transpose is invertible. Product of invertible matrices is invertible.

Let $\mathbf{A}$ and $\mb{A}_i$ be invertible matrices, $\mathbf{B}$ be a matrix of the same size of $\mathbf{A}$.
\[\mathbf{A}^{-1}=\frac{\mathrm{adj}\qty(\mathbf{A})}{\det\qty(\mathbf{A})}\]
\[\qty(\mathbf{A}^\top)^{-1}=\qty(\mathbf{A}^{-1})^\top\]
\bpr
\[\mb{A}^{\top}\qty(\mb{A}^{-1})^{\top}=\qty(\mb{A}^{-1}\mb{A})^{\top}=\mb{I}^\top=\mb{I}\]
\epr
\[\qty(\prod_{i=1}^n\mathbf{A}_i)^{-1}=\prod_{i=1}^n\qty(\mathbf{A}_{n-i+1})^{-1}\]
\[(\mathbf{A}^{-1}\mathbf{B}\mathbf{A})^n=\mathbf{A}^{-1}\mathbf{B}^n\mathbf{A},\quad n\in\mathbb{N}\]
Let $A,B$ be square matrix of same size,
\[AB=I\Rightarrow(\exists A^{-1}\land A^{-1}=B)\]
\subsubsection{Inverse of elementary matrix}
Every elementary matrix is invertible. Moreover, the inverse of an elementary is an elementary matrix.
\bit
\item Row addition: The inverse of $I_m$ with the $j,i$th entry replaced with $k$ is $I_m$ with the $j,i$th entry replaced with $-k$.
\item Row multiplication: The inverse of $I_m$ with the $i,i$th entry replaced with $k$ is $I_m$ with the $i,i$th entry replaced with $\frac{1}{k}$.
\item Row interchange: The inverse of $I_m$ with row $i$ and row $j$ interchanged is itself.
\eit
Thus, for any two row-equivalent $m\times n$ matrices $A,B$, there \forall n\in\bbZ\colonexists an $m\times m$ invertible matrix $C$, which is the product of a sequence of $m\times m$ elementary matrices, such that $B=CA$.
\subsection{Determinant}
\subsubsection{Determinant (行列式或判別式), cofactor (餘因子), cofactor matrix, and adjugate matrix (伴隨矩陣) or classical adjoint (經典伴隨矩陣)}
Let $\mathbf{A}_{ij}$ denote the submatrix of $\mathbf{A}$ without the $i$th row and $j$th column.

The determinant of a $1\times 1$ matrix is defined as the only element of it; the determinant of an $n\times n$ matrix $\mathbf{A}$, denoted $\det(\mathbf{A})$ or 
\[\begin{vmatrix}a_{11} & \ldots & a_{1n}\\\vdots & \ddots & \vdots\\a_{n1} & \ldots & a_{nn}\end{vmatrix},\]
is defined recursively as, for any $j\in\bbN_{\le n}$,
\[\det(\mathbf{A})=\sum_{i=1}^na_{ij}(-1)^{i+j}\det(\mathbf{A}_{ij})\]
or, for any $i\in\bbN_{\le n}$,
\[\det(\mathbf{A})=\sum_{j=1}^na_{ij}(-1)^{i+j}\det(\mathbf{A}_{ij})\]
The $(i,j)$-cofactor, denoted $c_{ij}$, of $\mb{A}$, is defined as
\[c_{ij}=(-1)^{i+j}\det(\mathbf{A}_{ij})\]
Thus,
\[\det(\mathbf{A})=\sum_{i=1}^na_{ij}c_{ij},\]
called the cofactor expansion of $\mb{A}$ along the column $j$; and
\[\det(\mathbf{A})=\sum_{j=1}^na_{ij}c_{ij},\]
called the cofactor expansion of $\mb{A}$ along the row $i$.

The cofactor matrix $\mathrm{cof}(\mathbf{A})$ of $\mathbf{A}$ is an $n\times n$ matrix defined as
\[\mathrm{cof}(\mathbf{A})_{ij}=c_{ij}\]
The adjugate matrix or classical adjoint matrix $\mathrm{adj}(\mathbf{A})$ of $\mathbf{A}$ is an $n\times n$ matrix defined as
\[\mathrm{adj}(\mathbf{A})=\mathrm{cof}(\mathbf{A})^\top\]
\subsubsection{Determinant identities}
Let $A$ and $A_i$ be square matrices.
\[\det\qty(\prod_{i=1}^nA_i)=\prod_{i=1}^n\det\qty(A_i)\]
\[\det(A)\neq 0\iff\exists A^{-1}\implies\det(A^{-1})=\frac{1}{\det(A)}\]
\[\det(\mb{I}_n)=1\]
\subsubsection{Determinant under elementary row operations}
For a square matrix $A$ and an elementary row operation that transforms $A$ to $B$,
\bit
\item Row addition $R_i\leftarrow R_i+cR_j$: $\det(B)=\det(A)$.
\item Row multiplication $R_i\leftarrow cR_i$: $\det(B)=c\det(A)$.
\item Row interchange $R_i\leftrightarrow R_j$: $\det(B)=\det(A)$.
\eit
\subsection{Special matrices}
\subsubsection{Diagonal matrix (對角矩陣)}
A diagonal matrix is a square matrix whose all off-diagonal entries are the additive identity of $0_{F}$.
\subsubsection{Diagonal matrix inverse}
Let
\[\mathbf{D} = \begin{bmatrix}
d_1 & 0 & \cdots & 0 \\
0 & d_2 & \cdots & 0 \\
\vdots & \vdots & \ddots & \vdots \\
0 & 0 & \cdots & d_n
\end{bmatrix}\]
where all $d_i\neq 0$. Then:
\[\mathbf{D}^{-1} = \begin{bmatrix}
\frac{1}{d_1} & 0 & \cdots & 0 \\
0 & \frac{1}{d_2} & \cdots & 0 \\
\vdots & \vdots & \ddots & \vdots \\
0 & 0 & \cdots & \frac{1}{d_n}\end{bmatrix}\]
\subsubsection{Diagonal matrices product}
Let $\mathbf{E}$, $\mathbf{F}$ be two $n\times n$ diagonal matrices. Then:
\[\mathbf{E}\mb{F} = \begin{bmatrix}
e_1f_1 & 0 & \cdots & 0 \\
0 & e_2f_2 & \cdots & 0 \\
\vdots & \vdots & \ddots & \vdots \\
0 & 0 & \cdots & e_nf_n
\end{bmatrix}\]
\subsubsection{Lower triangular matrix}
A lower triangular matrix is a matrix whose all entries $a_{ij}$ with $i\le j$ are the additive identity of $0_{F}$.
\subsubsection{Upper triangular matrix}
An upper triangular matrix is a matrix whose all entries $a_{ij}$ with $i\ge j$ are the additive identity of $0_{F}$.
\subsubsection{Lower triangular square matrices product}
Let $A,B$ be two $n\times n$ lower triangular matrices. Then $AB$ is an $n\times n$ lower triangular matrix.
\bpr
\[(AB)_{ij}=\sum_{k=1}^na_{ik}b_{kj}\]
For $i<j$,
\bit
\item If $k>i$, then $a_{ik}=0$.
\item If $k\le i$, then $k<j$, so $b_{kj}=0$
\eit
Thus $(AB)_{ij}=0$.
\epr
\subsubsection{Upper triangular square matrices product}
Let $A,B$ be two $n\times n$ upper triangular matrices. Then $AB$ is an $n\times n$ upper triangular matrix.
\bpr
\[(AB)_{ij}=\sum_{k=1}^na_{ik}b_{kj}\]
For $i>j$,
\bit
\item If $k<i$, then $a_{ik}=0$.
\item If $k\ge i$, then $k>j$, so $b_{kj}=0$
\eit
Thus $(AB)_{ij}=0$.
\epr
\sssc{Triangular square matrix determinant}
The determinant of a diagonal matrix, an upper triangular square matrix, or a lower triangular square matrix equals the product of its diagonal entries.
\bpr
Expand over a row or column with only one nonzero entry.
\epr
\subsubsection{Vandermonde matrix (范德蒙矩陣)}
A Vandermonde matrix $\mathbf{V}$ is:
\[V_{i,j}=\alpha_i^{\pht{i}j-1}\]
An $n\times n$ Vandermonde matrix $\mathbf{V}$ satisfies:
\[\det(V)=\prod _{1\leq i<j\leq n}(\alpha _{j}-\alpha _{i})\]
\ssc{Linear map and matrix}
\sssc{Linear map and matrix are isomorphic in finite-dimensional vector spaces with chosen ordered basis}
Let $V,W$ be finite-dimensional vector spaces of dimensions $m,n$ over a field $F$. Then, $L(V,W)$ and $F^{n\times m}$ are isomorphic.

Moreover, given arbitrary chosen ordered bases $B=(v_1,\ldots,v_m)$ for $V$ and $C=(w_1,\ldots,w_n)$ for $W$, for any $T\in L(V,W)$, each $T(v_j)$ can be represented as
\[T(v_j)=\sum_{i=1}^na_{ij}w_i\]
for some unique scalars $a_{ij}\in F$, and an isomorphism $f\colon L(V,W)\to F^{n\times m}$, which is a linear map, can be constructed as
\[f(T)=\begin{bmatrix}a_{11}&\cdots&a_{1m}\\\vdots&\ddots&\vdots\\a_{n1}&\cdots&a_{mn}\end{bmatrix}\]
and is such that
\[T(v)=f(T)v\]
for all $v\in V$, with vectors in $V$ in coordinate notation in terms of $B$ and vectors in $W$ in coordinate notation in terms of $C$.

And $f^{-1}$, which is a linear map, can be constructed as, for any $A=\begin{bmatrix}a_{11}&\cdots&a_{1m}\\\vdots&\ddots&\vdots\\a_{n1}&\cdots&a_{mn}\end{bmatrix}\in F^{n\times m}$,
\[(f^{-1}(A))(\sum_{j=1}^mx_jv_j)=\sum_{j=1}^mx_j\sum_{i=1}^na_{ij}w_i\]
and is such that
\[Av=(f^{-1}(A))(v)\]
for all $v\in V$, with vectors in $V$ in coordinate notation in terms of $B$ and vectors in $W$ in coordinate notation in terms of $C$.

Consequently, $L(V,W)$ is a vector space of dimension $mn$.
\bpr
Since $C$ is a basis, $a_{ij}$ are unique.

Write $v$ as a linear combination of $v_1,\ldots,v_m$. By linearity of $T$, we have
\[T(v)=f(T)v\]
By linearity of matrix multiplication, we have
\[Av=(f^{-1}(A))(v)\]
$f$ is injective: If $f(S)=f(T)$, since $a_{ij}$ are unique, $S(v_j)=T(v_j)$ for all $j$. Since $B$ spans $V$ and $S,T$ are linear, $S(v)=T(v)$ for all $v\in V$.

$f^{-1}$ as constructed is the inverse of $f$: Every $v\in V$ has a unique coordinate notation in terms of $B$, so $f^{-1}(A)$ is well-defined for all $A\in F^{n\times m}$. By construction, $f(f(^{-1})(A))=A$. Thus $f$ is surjective and $f^{-1}$ is the inverse of $f$.
\epr
$f(T)$ is called the standard matrix of $T$, and chosen ordered bases are often assumed.
\sssc{Composition of linear map}
For two linear maps $S\colon U\to V,T\colon V\to W$ with standard matrix $A,B$, its composition $T\circ S$, often denoted $TS$, has standard matrix $BA$.
\sssc{Inverse as linear map is inverse as matrix}
A linear map $T$ between two finite-dimensional vector spaces is bijective iff its standard matrix $A$ is injective, in which case the standard matrix of $T^{-1}$ is $A^{-1}$.
\sssc{Transpose as linear map is transpose as matrix}
For a linear map $T$ between two finite-dimensional vector spaces with standard matrix $A$, the standard matrix of the transpose (aka dual map or adjoint) $^tT$ is $A^{\top}$.
\subsection{System of Linear Equations (線性方程組)}
\subsubsection{System of equations (方程組)}
A solution to a system of equations is a choice of variables such that all equations are satisfied. Two systems of equations are called equivalent if they have exactly the same set of solutions.
\begin{itemize}
\item\tb{Consistent system (相容方程組)}: A system of equations is consistent iff it has at least one solution.
\item\tb{Inconsistent/contradictory system (相容方程組)}: A system of equations is inconsistent or contradictory iff it has no solution.
\item\tb{Dependent system (相依方程組)}: A consistent system of equations is dependent iff it has infinitely many solutions.
\item\tb{Independent system (獨立方程組)}: A system of equations is independent iff it has finitely many solutions.
\end{itemize}`
\subsubsection{Linear Equation (線性方程)}
A linear equation is an equation that may be put in the from:
\[\sum_{k=1}^na_kx_k=c,\]
where $a_k$ are called coefficients and $c$ is called constant term.

If $c=0$, the linear equation is called a homogeneous linear equation; otherwise, it is called a heterogeneous or non-homogeneous linear equation.
\subsubsection{System of Linear Equations (線性方程組) or linear system}
A system of linear equations is a collection of two or more linear equations involving the same tuple of variables.
\subsubsection{Independent system of linear equations has exactly one solution}
An independent system of linear equations must have exactly one solution.
\bpr
If more than one solution exist, their linear combinations form infinitely many solutions.
\epr
\subsubsection{Dependent system general solution}
For a dependent system, we can solve some of the variables, called basic variables, in terms of some other variables, called free variables, such that the free variables are independent from each other and that a choice of free variables uniquely determines basic variables, providing a general solution of the system.
\subsubsection{Coefficient matrix (係數矩陣)}
The coefficients of a system of linear equations written as a matrix, with one row for each equation.
\subsubsection{Constant term vector}
The column vector of the constant terms on the right side of the equal signs of a system of linear equations.
\subsubsection{Augmented matrix (增廣矩陣)}
Coefficient matrix with the constant term vector augmented to the last column of the matrix. For coefficient matrix $A$ and constant term vector $b$, we write the augmented matrix as $[A|b]$.
\subsubsection{Row equivalent means same solutions}
Two systems of linear equations with same size of augmented matrix have the same solutions if their augmented matrices are row-equivalent.

Two consistent systems of linear equations with same size of augmented matrix have the same solutions iff their augmented matrices are row-equivalent.
\bpr
Let $M=[A|b]$ and $N=[C|d]$ both be with $m$ rows and $n+1$ columns.

If: There exists elementary matrices $E_1,\ldots, E_k$ such that
\[N=E_k\ldots E_1M\]
We prove that multiplying augmented matrix $E$ by an elementary matrix doesn't change its solution set.
\[EM=[EA|Eb]\]
If $b\notin\row(A)$, since elementary row operation doesn't alter row space, $Eb\notin\row(EA)$. Otherwise, since $E$ is invertible with inverse being elementary matrix, we have
\[Ax=b\implies EAx=EAb\]
\[EAx=Eb\implies E^{-1}EAx=E^{-1}EAb\implies Ax=b\]

Only if: By the existence and uniqueness of reduced row echelon form, if they have the same non-empty solution set, they have the same reduced row echelon form $R$, and we can construct sequences of elementary matrices $E_1,\ldots,E_k$ and $F_1,\ldots,F_k$ such that
\[R=E_k\ldots E_1M=F_k\ldots F_1N\]
Since an elementary matrix is invertible with inverse being elementary matrix, we have
\[M=E_1^{-1}\ldots E_k^{-1}F_k\ldots F_1N,\]
and that $M,N$ are row-equivalent.
\epr
\subsubsection{Cramer Rule or Cramer formula (克拉瑪公式)}
Given a system of linear equations that can be written in the form:
\[\mathbf{A}\mathbf{X}=\mathbf{C}\]
where $\mathbf{A}$ is an $n\times n$ coefficient matrix, $\mathbf{X}$ is a variable matrix:
\[\mathbf{X}=\begin{pmatrix}x_1, x_2, x_3, \ldots, x_n\end{pmatrix}^\top,\]
and $\mathbf{C}$ is a coefficient matrix:
\[\mathbf{C}=\begin{pmatrix}c_1, c_2, c_3, \ldots, c_n\end{pmatrix}^\top.\]
Let \(\mathbf{A}_i\) be $\mathbf{A}$ with the \(i\)th column replaced by \(\mathbf{C}\). The solution to the system of linear equations is given by:
\[\begin{cases}
x_i=\frac{\det\qty(\mathbf{A}_i)}{\det\qty(\mathbf{A})},\quad\det\qty(\mathbf{A})\neq 0\\
\text{inconsistent},\quad\det\qty(\mathbf{A})=0\land\bigvee_{i=1}^n\det\qty(\mathbf{A}_i)\neq 0\\
\text{dependent},\quad\det\qty(\mathbf{A})=0\land\bigvee_{i=1}^n\det\qty(\mathbf{A}_i)=0
\end{cases}\]
where $x_i=\frac{\det\qty(\mathbf{A}_i)}{\det\qty(\mathbf{A})}$ can be geometrically viewed as the volume of the parallelogram formed by the column vectors of $\mathbf{A}_i$ divided by the volume of the parallelogram formed by the column vectors of $\mathbf{A}$.

Particularly, when $\mathbf{C}$ is a zero matrix, that is,
\[\mathbf{A}\mathbf{X}=\mathbf{0},\]
the solution to the system of linear equations is given by:
\[\begin{cases}
x_i=0,\quad\det\qty(\mathbf{A})\neq 0\\
\tx{dependent},\quad\det\qty(\mathbf{A})=0
\end{cases}\]
\subsection{Row echelon form}
\subsubsection{(General) row echelon form (REF) (（廣義）zh-hant: 列/zh-hans: 行階梯形式)}
A matrix is said to be in row echelon form iff:
\ben
\item Each nonzero lies above every zero row.
\item The leading entry of every nonzero row lies in a column to the right of the column containing the leading entry of any row above.
\een
Thus, for a leading entry of any row, all entries below it in the same column are zero.
\subsubsection{Reduced row echelon form (RREF)}
A matrix is said to be in reduced row echelon form iff:
\ben
\item It is in row echelon form.
\item Each leading entry is the only nonzero entry in the column.
\item The leading entries of all nonzero rows are 1.
\een
\subsubsection{Pivot position}
For a matrix $A$ and its (unique, which will be proven later) reduced row echelon form $R$, the position containing a leading entry of a nonzero row is called pivot position, a column containing a pivot position is called pivot column, and a row containing a pivot position is called pivot row.
\subsubsection{Gaussian Elimination (高斯消去法), row reduction, Gauss-Jordan Elimination, and reduced row echelon form test for consistency and obtaining general solution}
Gaussian elimination or row reduction refers to either algorithm, or the forward pass only with Gaussian-Jordan elimination referring to either the entire algorithm or the backward pass only.

To find the row echelon form of a matrix, do the followings, called forward pass:
\ben
\item Determine the leftmost nonzero column. This is a pivot column, and the topmost position in this column is a pivot position.
\item In the pivot column, choose any nonzero entry, and perform the appropriate row interchange to bring this entry into the pivot position.
\item Add an appropriate multiple of the row containing the pivot position to each lower row in order to change each entry below the pivot position into 0.
\item Ignore the row containing the pivot position. If there is a nonzero row below it, repeat steps 1–4 on the submatrix that remains.
\een
To find the reduced row echelon form of a matrix from a row echelon form of it, from the topmost to bottommost nonzero row, do the followings, called backward pass:
\ben
\item Perform the appropriate row scaling to make the leading entry 1.
\item Add an appropriate multiple of the row to each upper row in order to change each entry above the leading entry into 0.
\een
A system of linear equations is consistent iff the reduced row echelon form of its augmented matrix has no row with leading entry in the rightmost column.
\bpr
If the reduced row echelon form of its augmented matrix has row with leading entry in the rightmost column, it is obviously inconsistent. If the reduced row echelon form of its augmented matrix has no row with leading entry in the rightmost column, we can construct a general solution.
\epr
Given the reduced row echelon form of the augmented matrix of a linear system, if it has no row with leading entry in the rightmost column, we can obtain the general solution by
\ben
\item Assign the variables of non-pivot columns as free variables and variables of pivot columns as basic variables.
\item For each nonzero row, write an equation of the variable of the leading entry as a linear combination of the variables (constant term for the rightmost column and free variables for others) of other nonzero entries.
\een
\subsection{Subspaces of matrices and linear maps}
\subsubsection{Four fundamental subspaces of a matrix (矩陣的四大子空間)}
The four fundamental subspaces of a matrix are the row space, column space, null space, and left null space.
\subsubsection{Column space (zh-hant: 行/zh-hans: 列空間)}
The column space of $\mathbf{A}$, denoted as \(\operatorname{Col}(\mathbf{A})\) or $\mathrm{\mathbf{C}}(\mathbf{A})$, is defined as the span of the set of column vectors of $\mathbf{A}$.
\subsubsection{Row space (zh-hant: 列/zh-hans: 行空間)}
The row space of $\mathbf{A}$, denoted as \(\operatorname{Row}(\mathbf{A})\), is defined as the span of the set of the row vectors of $\mathbf{A}$, that is, $\operatorname{Col}(\mathbf{A}^{\top})$.
\subsubsection{Null space (零空間), kernel (核), or kernel space (核空間)}
The kernel of a linear map $L\colon V\to W$, denoted as $\ker(L)$ or $\Null(L)$, is defined as:
\[\ker(L)=\{v\in V:\,L(v)=0\}.\]
$\dim\ker(\mathbf{A})$ is called nullity (零化度).
\subsubsection{Left null space (左零空間), left kernel (左核), or left kernel space (左核空間)}
The left null space of a linear map $L\colon V\to W$ is defined as $\ker(L^{-1})$.

$\dim\ker(L^{-1})$ is called left nullity.
\subsubsection{Image space (像空間) or range (值域)}
The image space of a linear map $T\colon V\to W$, denoted as $\operatorname{im}(T)$, is defined as,
\[\operatorname{im}(T)=\{T(v):\,v\in V\}.\]
\subsubsection{Cokernel (餘核)}
The cokernel of a linear map $T\colon V\to W$, denoted as $\operatorname{coker}(T)$, is defined as,
\[\operatorname{coker}(T)=W\setminus\operatorname{im}(T).\]
\subsection{Rank and nullity}
\subsubsection{Rank (秩)}
The rank of a matrix \(\mathbf{A}\), denoted as \( \operatorname{rank}(\mathbf{A}) \) or $\operatorname{rk}(\mathbf{A})$, is defined as the dimension of its column space (or row space, which are the same by constant rank theorem), that is,
\[\text{rank}(\mathbf{A})=\dim\operatorname{Col}(\mathbf{A})\]
\subsubsection{Constant rank theorem (常秩定理)}
For any matrix $\mathbf{A}$:
\[\dim\operatorname{Col}(\mathbf{A})=\dim\operatorname{Row}(\mathbf{A})\]
\begin{proof}
Let $\mathbf{A}$ be an $m\times n$ matrix, $r=\dim\operatorname{Col}\qty(\mathbf{A})$, $ c_{1},c_{2},\ldots ,c_{r}$ be a basis of $\operatorname{Col}(\mathbf{A})$, matrix $\mathbf{C}=[c_{1},c_{2},\ldots ,c_{r}]$, and $r\times n$ matrix $\mathbf{R}$ such that $\mathbf{A}=\mathbf{C}\mathbf{R}$. 

Since:
\[\operatorname{Row}(\mathbf{A})\subseteq\operatorname{Row}(\mathbf{R})\]
\[\dim\operatorname{Row}(\mathbf{A})\leq\dim\operatorname{Row}(\mathbf{R})\leq r=\dim\operatorname{Col}(\mathbf{A})\]
Similarly, we can obtain that
\[\dim\operatorname{Col}(\mathbf{A})\leq\dim\operatorname{Row}(\mathbf{A})\]
Therefore,
\[\dim\operatorname{Col}(\mathbf{A})=\dim\operatorname{Row}(\mathbf{A})\]
\end{proof}
\subsubsection{Rank-nullity theorem (秩—零化度定理)}
For any $m\times n$ matrix $\mathbf{A}$:
\[\rank(\mathbf{A})+\dim(\ker(\mathbf{A}))=n\]
\begin{proof}
Let \(\mathbf{A}\) be an \(m\times n\) matrix with rank \(r\). We will show that there exists a set of \(n-r\) linearly independent solutions to the homogeneous system \(\mathbf{Ax}=\mathbf{0}_m\) where $\mb{x}$ is an $n$-dimensional column vector.

To do this, we will produce an \(n\times(n-r)\) matrix \(\mathbf{X}\) whose columns form a basis of the null space of \(\mathbf{A}\).

Without loss of generality, assume that the first \(r\) columns of \(\mathbf{A}\) are linearly independent. So, we can write
\[\mathbf{A}=\begin{pmatrix}\mathbf{A}_1&\mathbf{A}_2\end{pmatrix},\]
where
\begin{itemize}
\item \(\mathbf{A}_1\) is an \(m \times r\) matrix with \(r\) linearly independent column vectors, and
\item \(\mathbf{A}_2\) is an \(m\times(n-r)\) matrix such that each of its \(n-r\) columns is a linear combination of the columns of \(\mathbf{A}_1\).
\end{itemize}
This means that \(\mathbf{A}_2 = \mathbf{A}_1\mathbf{B}\) for some \(r \times (n-r)\) matrix \(\mathbf{B}\), and hence,
\[\mathbf{A} = \begin{pmatrix} \mathbf{A}_1 & \mathbf{A}_1\mathbf{B}\end{pmatrix}.\]
Let
\[\mathbf{X} = \begin{pmatrix} -\mathbf{B} \\ \mathbf{I}_{n-r} \end{pmatrix},\]
Then,
\[\mathbf{A}\mathbf{X} = \begin{pmatrix}\mathbf{A}_1 & \mathbf{A}_1\mathbf{B} \end{pmatrix}\begin{pmatrix} -\mathbf{B} \\ \mathbf{I}_{n-r} \end{pmatrix} = -\mathbf{A}_1\mathbf{B} + \mathbf{A}_1\mathbf{B} = \mathbf{0}_{m \times (n-r)}.\]
Therefore, each of the \(n-r\) columns of \(\mathbf{X}\) is a solution of \(\mathbf{Ax} = \mathbf{0}_m\).

Furthermore, we claim that the \(n-r\) columns of \(\mathbf{X}\) are linearly independent. We show that $\mathbf{0}_{n-r}$ is the only $(n-r)$-dimensional column vector $\mb{u}$ such that
\[\mathbf{Xu}=\mathbf{0}_n\]
\[\begin{pmatrix}
-\mathbf{B} \\
 \mathbf{I}_{n-r}
\end{pmatrix}\mathbf{u} = \begin{pmatrix}
-\mathbf{B}\mathbf{u} \\
 \mathbf{u}
\end{pmatrix} = \begin{pmatrix}
\mathbf{0}_r \\
 \mathbf{0}_{n-r}
\end{pmatrix}\]
We next prove that any solution of \(\mathbf{Ax} = \mathbf{0}_{m}\) must be in the column space of \(\mathbf{X}\). Let the submatrix of the top $r$ rows of $\mb{x}$ be $\mb{x}_1$ and the submatrix of the bottom $n-r$ rows of $\mb{x}$ be $\mb{x}_2$.
\[\mathbf{A}\mathbf{x} = \begin{pmatrix}\mathbf{A}_1 & \mathbf{A}_1\mathbf{B}\end{pmatrix} \begin{pmatrix} \mathbf{x}_1 \\ \mathbf{x}_2 \end{pmatrix} = \mathbf{A}_1\mathbf{x}_1 + \mathbf{A}_1\mathbf{B}\mathbf{x}_2 = \mathbf{A}_1(\mathbf{x}_1 + \mathbf{B}\mathbf{x}_2) = \mathbf{0}_{m} \]
Since the columns of \(\mathbf{A}_1\) are linearly independent, \(\mathbf{A}_1\mathbf{u} = \mathbf{0}_{r}\) implies \(\mathbf{u} = \mathbf{0}_{r}\), thus
\[\mathbf{u}_1 + \mathbf{B}\mathbf{u}_2  = \mathbf{0}_{r} \]
\[\mathbf{u} = \begin{pmatrix} -\mathbf{B} \\ \mathbf{I}_{n-r} \end{pmatrix}\mathbf{u}_2 = \mathbf{X}\mathbf{u}_2\]

This proves that any vector \(\mathbf{u}\) that is a solution of \(\mathbf{Ax} = \mathbf{0}\) must be a linear combination of the \(n-r\) special solutions given by the columns of \(\mathbf{X}\). And we have already seen that the columns of \(\mathbf{X}\) are linearly independent. 

Hence, the column space of \(\mathbf{X}\) is the null space of \(\mathbf{A}\). Therefore, the nullity of \(\mathbf{A}\) is \(n - r\). Since \(r\) is the rank of \(\mathbf{A}\), it follows that \(\text{rank}(\mathbf{A}) + \ker(\mathbf{A}) = n\).
\end{proof}
\subsubsection{Rank-left nullity theorem}
For any $m\times n$ matrix $\mathbf{A}$:
\[\rank(\mathbf{A})+\dim(\ker(\mathbf{A}^{\top}))=m\]
\bpr
Apply rank-nullity theorem on $A^{\top}$.
\epr
\subsubsection{Full rank (滿秩) and rank deficiency (秩虧)}
For any $m\times n$ matrix $\mb{A}$:
\[\text{rank}(\mathbf{A}) \leq \min(m, n).\]
$\mathbf{A}$ is called to have full rank iff
\[\text{rank}(\mathbf{A}) = \min(m, n)\]
otherwise it is called to be rank-deficient (欠秩或秩虧的).

The rank deficiency of a $m\times n$ matrix $\mb{A}$ is defined as
\[\min(m, n)-\text{rank}(\mathbf{A}) \]
\subsubsection{Zero rank}
A matrix is a zero matrix iff its rank is zero.
\subsubsection{Rank preservation}
For real matrix,
\[\rank(A)=\rank(A^{\top})=\rank(AA^{\top})=\rank(A^{\top}A)\]
For complex matrix,
\[\rank(A)=\rank(A^{\top})=\rank(\bar{A})=\rank(A^*)=\rank(AA^*)=\rank(A^*A)\]
\subsection{Properties}
\subsubsection{Column correspondence property}
Let $A,B$ be row-equivalent matrices. If a column $j$ of $A$ is a linear combination of other columns of $A$, then column $j$ of $B$ is a linear combination of the corresponding columns of $B$ using the same coefficients.
\bpr
There exists an invertible matrix $E$, which is the product of a sequence of $m\times m$ elementary matrices, such that $B=EA$. Let a matrix with subscript denote its submatrix of that column. For every column $j$
\[B_j=EA_j\]
Suppose
\[A_j=\sum_{i\ne j}c_iA_i\]
\[B_j=EA_j=E\sum_{i\ne j}c_iA_i=\sum_{i\ne j}c_iB_i\]
\epr
\subsubsection{Elementary row operation preservations}
Suppose $E$ is an $m\times m$ elementary matrix and $A$ is an $m\times n$ matrix. Then $A$ and $B\coloneq EA$ have the same
\bit
\item row space (and thus also rank)
\item kernel (and thus also nullity)
\eit
\bpr
Since $B=EA$, each row of $B$ is a linear combination of rows of $A$, thus $\Row(B)\subseteq\Row(A)$. Since $A=E^{-1}B$, each row of $A$ is a linear combination of rows of $B$, thus $\Row(A)\subseteq\Row(B)$. Thus $\Row(A)=\Row(B)$.

$Bx=0\iff EAx=0\iff Ax=E^{-1}0=0\iff EAx=Bx=E0=0$, thus $\ker(A)=\ker(B)$.
\epr
The statement does not hold for column space and left kernel.
\bpr
Counterexample for column space:
\[A=\begin{bmatrix}1\\0\end{bmatrix}\]
\[B=\begin{bmatrix}0\\1\end{bmatrix}\]
$B$ is obtained by row swapping $A$, but $\Col(A)=\span\qty{\begin{bmatrix}1\\0\end{bmatrix}}$ and $\Col(B)=\span\qty{\begin{bmatrix}0\\1\end{bmatrix}}$ are different.

Counterexample for left kernel:
\[A=\begin{bmatrix}1&0\\0&0\end{bmatrix}\]
\[B=\begin{bmatrix}1&0\\1&0\end{bmatrix}\]
$B$ is obtained by adding row 1 of $A$ to row 2, but $\ker(A^\top)=\span\qty{\begin{bmatrix}0\\1\end{bmatrix}}$ and $\ker(B^\top)=\span\qty{\begin{bmatrix}1\\-1\end{bmatrix}}$ are different.
\epr
\subsubsection{Row-equivalence is an equivalence relation}
Row-equivalence is an equivalence relation.
\bpr
Reflexivity: Obviously

Symmetry: Since elementary matrices are invertible

Transitivity: If $B$ is row-equivalent to $A$ with $B=EA$ and $C$ is row-equivalent to $B$ with $C=FB$, then $C=EFA$, thus $C$ is row-equivalent to $A$.
\epr
\subsubsection{Rank and number of pivot columns}
For a matrix $A$ and its reduced row echelon form $R$, the followings are equal:
\bit
\item The rank of $A$ (or $R$)
\item The number of pivot positions of $A$ (or $R$)
\item The number of pivot columns of $A$ (or $R$)
\item The number of pivot rows of $A$ (or $R$)
\item The number of nonzero rows of $R$
\item The number of basic variables in the general solution of any consistent system of linear equations whose coefficient matrix is $A$ (or $R$)
\eit
\subsubsection{Nullity and number of non-pivot columns}
For a matrix $A$ and its reduced row echelon form $R$, the followings are equal:
\bit
\item The nullity of $A$ (or $R$)
\item The number of non-pivot columns of $A$ (or $R$)
\item The number of free variables in the general solution of any consistent system of linear equations whose coefficient matrix is $A$ (or $R$)
\eit
\subsubsection{Left nullity and number of non-pivot rows}
For a matrix $A$ and its reduced row echelon form $R$, the followings are equal:
\bit
\item The left nullity of $A$ (or $R$)
\item The number of non-pivot rows of $A$ (or $R$)
\item The number of zero rows of $R$
\item The number of free variables in the general solution of any consistent system of linear equations whose coefficient matrix is $A^\top$ (or $R^{\top}$)
\eit
\subsubsection{Existence and uniqueness of reduced row echelon form}
Every matrix can be transformed into one and only one matrix in reduced row echelon form by means of a sequence of elementary row operations.
\bpr
Existence: By Gauss-Jordan elimination.

Uniqueness:

Let $A$ be a matrix and $R$ be any reduced row-equivalent form of it. Since elementary row operation preserves rank, which is the dimension of column space, which is the number of pivot columns of both $A$ and $R$, the pivot columns of $R$ are uniquely determined.

Let a matrix with subscript denote its submatrix of that column. For every non-pivot column $j$, by column correspondence property, since column $j$ of $A$ is a linear combination of pivot columns of $A$, column $j$ of $R$ is a linear combination of the corresponding columns of $A$ using the same coefficients, which is unique.
\epr
\subsubsection{Column space test for consistency of linear system}
The linear system $Ax=b$ is consistent iff $b\in\Col(A)$.
\bpr
By definition, obviously.
\epr
Thus, a homogeneous linear system is always consistent.
\bpr
Zero vector is in all linear subspaces.
\epr
\subsubsection{Invertible coefficient matrix}
For a linear system $Ax=b$ with invertible coefficient matrix $A$ and constant term vector $b$, then $x=A^{-1}b$ is the unique solution to it.
\subsubsection{Solution space is an affine space with kernel as associated vector space}
Suppose $y$ is a solution to a linear system $Ax=b$ with coefficient matrix $A$ and constant term vector $b$, then the set of all solutions of it
\[\{y+k\colon k\in\ker(A)\}\]
\bpr
Suppose $z-y\in\ker(A)$.
\[Az=Ay+A(z-y)=Ay=b\]
Suppose $z-y\notin\ker(A)$.
\[Az=Ay+A(z-y)=b+A(z-y)\ne b\]
\epr
Thus, the solution space is an affine space with the kernel of its coefficient matrix as the associated vector space. For homogeneous linear system, the solution space is the kernel of its coefficient matrix, whicj a vector space.
\subsubsection{Zero nullity}
For a matrix $A\in\bbR^{m\times n}$, the followings are equivalent:
\bit
\item $A$ has zero nullity (i.e., $\rank(A)=\dim\Col(A)=\dim\Row(A)=n$).
\item The linear system $Ax=b$ has at most one solution.
\bpr
If consistent, since the associated vector space $\ker(A)$ of solution space is $\{\mb{0}\}$, the solution space has only one element.
\epr
\item The only solution to the homogeneous linear system whose coefficient matrix is $A$ is $\mb{0}$.
\item The columns of the reduced row echelon form of $A$ are the $n$ distinct standard basis vectors of $F^n$.
\item The linear map whose standard matrix is $A$ is injective.
\eit
\subsubsection{Zero left nullity}
For a matrix $A\in\bbR^{m\times n}$, the followings are equivalent:
\bit
\item $A$ has zero left nullity (i.e., $\rank(A)=\dim\Col(A)=\dim\Row(A)=m$).
\item The linear system $Ax=b$ is consistent for all $b$.
\bpr
Since $\dim\Col(A)=m$, the column space is $\bbR^m$, thus the constant term vector is in the column space. By column space test for consistency of linear system, the system is consistent.
\epr
\item The linear map whose standard matrix is $A$ is surjective.
\eit
\subsubsection{Pivot column set is column space basis}
The set of pivot columns of a matrix is a basis of its column space.
\bpr
For reduced row echelon form, obviously. By column correspondence property and unique existence of reduced row echelon form, the statement holds.
\epr
\subsubsection{Invertible iff full rank}
A square matrix is invertible iff it has full rank.
\bpr
Implies: If $A$ is invertible $Ax=0$ has the only solution $x=0$, thus $A$ has nullity $0$. By rank-nullity thereom, $A$ has rank $n$.

Implied by: Since $A$ has full rank, elementary row operations preserve rank, and $I_n$ is the only RREF with full rank, there exists an invertible matrix $E$, which is the product of a sequence of elementary matrices, such that $EA=I_n$, then $E=A^{-1}$.
\epr
\subsubsection{Invertible matrix reduced row echelon form is identity matrix}
Let $A$ be an $n\times n$ matrix. Then $A$ is invertible iff its reduced row echelon form is $I_n$.
\bpr
Implies: Since $A$ has full rank, elementary row operations preserve rank, and $I_n$ is the only RREF with full rank, $A$ must have reduced row echelon form $I_n$.

Implied by: Since elementary row operations preserve rank and $I_n$ has full rank, $A$ has full rank, thus $A$ is invertible.
\epr
\subsubsection{Invertible matrix must be product of elementary matrix}
Let $A$ be an $n\times n$ matrix. Then $A$ is invertible iff it's a product of $n\times n$ elementary matrices.
\subsubsection{Algorithm for computing inverse with elementary row operations}
Let $A$ be an $n\times n$ matrix. Use elementary row operations to transform $[A I_n]$ into the form $[R B]$ where $R$ is the reduced row echelon form of $A$. Then either
\bit
\item $R=I_n$, in which case $A$ is invertible and $B=A^{-1}$; or
\item $R\ne I_n$, in which case $A$ is not invertible.
\eit
\subsubsection{Algorithm for computing inverse right multiplied by another matrix with elementary row operations}
Let $A$ be an $n\times n$ matrix and $B$ be an $n\times p$ matrix. Use elementary row operations to transform $[A B]$ into the form $[R C]$ where $R$ is the reduced row echelon form of $A$. Then either
\bit
\item $R=I_n$, in which case $A$ is invertible and $C=A^{-1}$; or
\item $R\ne I_n$, in which case $A$ is not invertible.
\eit
\subsection{Partitioning}
\subsubsection{Partitioning}
We can break up any matrix by drawing horizontal and vertical lines within the matrix, which divides the matrix into a two-dimensional array of submatrices called blocks. The resulting array is called a partition of the matrix, and the process of forming these blocks is called partitioning.
\subsubsection{Block multiplication}
Suppose two matrices $A$ and $B$ are partitioned into blocks so that the number of blocks in each row of $A$ is the same as the number of blocks in each column of $B$. Then the matrices can be multiplied according to the usual rules for matrix multiplication, treating the blocks as if they were scalars, provided that the individual products are defined.
\subsubsection{Matrix multiplication by rows}
Let $A$ be an $m\times n$ matrix and $B$ an $n\times p$ matrix. The $i$th row of $AB$ can be obtained by multiplying the $i$th row of $A$ by $B$.
\subsubsection{Matrix multiplication by outer product}
Let $A$ be an $m\times n$ matrix and $B$ an $n\times p$ matrix. The matrix $AB$ is the sum of matrix products of each column of $A$ with the corresponding row of $B$.
\subsection{Lower-upper (LU) decomposition or factorization}
\subsubsection{Lower-upper (LU) decomposition or factorization with partial pivoting}
LU decomposition with partial pivoting expresses an $m\times n$ matrix $A$ as the product a $m\times m$ lower triangular matrix $L$ whose all diagonal entries being multiplicative identity $1_F$ and an $m\times n$ upper triangular matrix $U$:
\[A=LU\]
If Gaussian elimination forward pass on $A$ doesn't need row interchange (i.e., only needs row additions). LU decomposition of $A$ exists. To obtain it, let $L_0$ be $I_m$, and perform Gaussian elimination forward pass on $A$ to obtain a row echelon form $U$ of $A$ with $l$ row additions. For each $k$th operation of adding row $i$ multiplied by $c$ to row $j$ in the process, assign $L_{k-1}$ with the $j,i$th entry replaced with $-c$ to $L_k$. Let $L=L_l$. We obtain the LU decomposition of $A$:
\[A=LU\]
\bpr
Observe that $U$ is an $m\times n$ upper triangular matrix. Let the elementary matrices of the operations be $E_1,\ldots,E_l$.
\[U=E_l\ldots E_1A\]
Since all $E_i$ are invertible,
\[A=E_1^{-1}\ldots E_l^{-1}U\]
Since all $E_i$ are lower triangular with all diagonal entries being multiplicative identity $1_F$,
\[L=E_1^{-1}\ldots E_l^{-1}\]
is lower triangular with all diagonal entries being multiplicative identity $1_F$.

In forward pass, no two same row additions are performed, thus
\[L_k=E_1^{-1}\ldots E_k^{-1}\]
\epr
\subsubsection{Lower-upper (LU) decomposition or factorization with partial pivoting}
LU decomposition with partial pivoting expresses an $m\times n$ matrix $A$ multiplied by an $m\times m$ permutation matrix $P$ whose columns and rows are standard basis vectors (i.e., each column and each row has exactly one multiplicative identity $1_F$ and $m-1$ additive identity $0_F$), as the product a $m\times m$ lower triangular matrix $L$ whose all diagonal entries being multiplicative identity $1_F$ and an $m\times n$ upper triangular matrix $U$:
\[PA=LU\]
Let $L_0=O_m$ and $P_0=I_m$. Perform Gaussian elimination forward pass on $A$ to obtain a row echelon form $U$ of $A$ with $l$ row additions and interchanges. For each elementary row operation $k$ in it,
\bit
\item if it is row addition that adds row $i$ multiplied by $c$ to row $j$, assign $L_{k-1}$ with the $j,i$th entry replaced with $-c$ to $L_k$ and assign $P_{k-1}$ to $P_k$;
\item if it is row interchange that swaps row $i$ and $j$, do the same operation on $L_{k-1}$ and $P_{k-1}$ to obtain $L_k$ and $P_k$.
\eit
Let $L=I_m+L_l$ and $P=P_l$. We obtain the LU decomposition with partial pivoting of $A$:
\[PA=LU\]
\bpr
Observe that $LU$ is the LU decomposition of $PA$, and observe the relation of the forward pass of $A$ and $PA$ to construct the algorithm for sequence $L_k$.
\epr
\subsubsection{Solving a system of linear equations using an LU decomposition}
The system $Ax=b$ can be rewritten as $LUx=b$. Set $y=Ux$.

First, solve $Ly=b$, called forward substitution. Let $y=(y_1,\ldots,y_m)^\top$ and $b=(b_1,\ldots,b_m)^\top$. If $m>n$, then
\bit
\item if $b_{n+1},\ldots,b_m$ are all zero, then the system is consistent, and we can discard the bottom $m-n$ rows,
\item otherwise, the system is inconsistent, and there is no solution.
\eit
Solve row by row from the top. The top row gives $y_1=b_1$. For each following row $i$, substituting results of rows above gives $y_i$.

Second, solve $Ux=y$, called backward substitution. Let $x=(x_1,\ldots,x_n)^\top$. If $m<n$, $x_{m+1},\ldots,x_n$ are free variables. Solve row by row from the bottom. The bottom row gives $x_m$ as a function of the free variables. For each following row $i$, substituting results of rows below gives $x_i$ as a function of the free variables.
\subsubsection{Solving a system of linear equations using an LU decomposition with partial pivot}
The system $Ax=b$ can be rewritten as $LUx=Pb$. Solve it using the LU decomposition $LU$.
\subsection{Flop count of various procedure}
Each floating point operation (flop) is an arithmetic operation of floating numbers. Suppose we have a system $Ax=b$ of floating numbers with $A$ being an $n\times n$ invertible matrix. The approximate flop counts are:
\bit
\item computing the LU decomposition of $A$: $\frac{2}{3}n^3$.
\item solving the system using Gaussian elimination: $\frac{2}{3}n^3$.
\item solving the system given the LU decomposition of $A$: $2n^2$.
\item calculating the inverse of $A$: $2n^3$
\eit
\bpr
TODO
\epr
\section{Matrix Analysis (矩陣分析)}
\subsection{Eigenvector or characteristic vector (特徵向量) and characteristic polynomial (特徵多項式)}
\subsubsection{Eigenvector or characteristic vector (特徵向量) and eigenvalue (本徵值、固有值、特徵值或特徵根)}
Given a linear transformation $L\colon V\to V$, an eigenvector of it is defined as any nonzero vector $v\in V$ such that:
\[\mathbf{A}v = \lambda v,\]
where $\lambda$ is a scalar, which is called the eigenvalue of $\mb{A}$ or the eigenvalue of $v$.

Obviously, $v$ is an eigenvector implies $av$ is an eigenvector of the same eigenvalue for any nonzero scalar $a$.
\subsubsection{Characteristic polynomial}
For an $n\times n$ matrix $\mathbf{A}$, the characteristic polynomial of it is defined as
\[p_{\mathbf{A}}(\lambda):=\det(\lambda \mathbf{I}_{n}-\mathbf{A}),\]
where $\lambda$ is a scalar variable.

A scalar $\lambda$ is an eigenvalue of $\mb{A}$ if and only if
\[p_\mathbf{A}(\lambda)=0.\]
\subsubsection{Multiplicity of eigenvalues}
The algebraic multiplicity of a root $\lambda$ of the characteristic polynomial of a square matrix is greater than or equal to the number of linearly independent eigenvectors of it whose eigenvalue is $\lambda$.
\bpr
TODO
\epr
\subsubsection{Number of eigenvalues}
The number of distinct eigenvalues of a square matrix is less than or equal to the number of linearly independent eigenvectors, which less than or equal to $n$.
\bpr
TODO
\epr
\subsubsection{Invertible iff has nonzero eigenvalue}
A square matrix is invertible iff it has nonzero eigenvalue.
\bpr
TODO
\epr
\subsubsection{Cayley–Hamilton theorem (凱萊–哈密頓定理)}
For a square matrix $\mathbf{A}$ and its characteristic polynomial $p_\mathbf{A}$, the Cayley–Hamilton theorem states that:
\[p_\mathbf{A}(\mathbf{A})=0\]
\subsubsection{Diagonalization (對角化)}
The diagonalization of an $n\times n$ matrix $\mb{A}$ is the construction of an $n\times n$ matrix $\mb{P}$ and an $n\times n$ diagonal matrix $\mb{D}$ such that:
\[\mathbf{A}=\mathbf{P}^{-1}\mathbf{D}\mathbf{P}\]

An $n\times n$ matrix can be diagonalized if it has $n$ distinct eigenvalues.

The steps of the diagonalization of an $n\times n$ matrix $\mathbf{A}$ is as follows:
\begin{enumerate}
\item For each eigenvalue $\lambda$ of the matrix $\mathbf{A}$, solve an eigenvector $v$:
\[(\mathbf{A} - \lambda \mathbf{I})\mathbf{v} = 0.\]
\item Compose matrix \( \mathbf{P} \) with each eigenvector of $\mb{A}$ as a column (or row).
\end{enumerate}

Common application:
\[\mathbf{A}^n=(\mathbf{P}^{-1}\mathbf{D}\mathbf{P})^n=\mathbf{P}^{-1}\mathbf{D}^n\mathbf{P}\]
\subsubsection{Definiteness (正定性)}
Let $A\in\bbR^{n\times n}$. $A$ is called
\bit
\item Positive definite: iff
\[\forall\mb{x}\in\bbR^n,\mb{x}\ne 0\colon\mb{x}^\top A\mb{x}>0.\]
\item Negative definite: iff
\[\forall\mb{x}\in\bbR^n,\mb{x}\ne 0\colon\mb{x}^\top A\mb{x}<0.\]
\item Positive semi-definite: iff
\[\forall\mb{x}\in\bbR^n,\mb{x}\ne 0\colon\mb{x}^\top A\mb{x}\geq 0.\]
\item Negative semi-definite: iff
\[\forall\mb{x}\in\bbR^n,\mb{x}\ne 0\colon\mb{x}^\top A\mb{x}\leq 0.\]
\item Indefinite: otherwise, that is,
\[\exists\mb{x}\in\bbR^n,\mb{x}\ne 0\colon\mb{x}^\top A\mb{x}>0\]
and
\[\exists\mb{x}\in\bbR^n,\mb{x}\ne 0\colon\mb{x}^\top A\mb{x}<0.\]
\eit
\subsubsection{Eigenvalues Test of Definiteness for Symmetric Hessians of Scalar Fields}
For a symmetric matrix $A$,
\bit
\item $A$ is positive definte if all eigenvalues of it are positive.
\item $A$ is negative definte if all eigenvalues of it are negative.
\item $A$ is positive semidefinte if all eigenvalues of it are nonnegative and at least one eigenvalue of it is positive.
\item $A$ is negative semidefinte if all eigenvalues of it are nonpositive and at least one eigenvalue of it is negative.
\item $A$ is neutral semidefinte if all eigenvalues of it are zero.
\item $A$ is indefinte otherwise.
\eit
\subsection{Generalized Definition of Inverse}
\subsubsection{One-sided inverse (單邊逆)}
\begin{itemize}
\item Right inverse (右逆): For any $n\times m$ matrix $A$ with rank $\mathrm {rank}(A)=n$, there exists an $m\times n$ matrix $A_{\mathrm {R} }^{-1}$, called the right inverse of $\mb{A}$, such that
\[AA_{\mathrm {R} }^{-1}=I_{n},\]
where $I_{n}$ is the $n\times n$ identity matrix.
\item Left inverse (左逆): For any $n\times m$ matrix $A$ with rank $\mathrm {rank}(A)=m$, there exists an $m\times n$ matrix $A_{\mathrm {L} }^{-1}$, called the left inverse of $\mb{A}$, such that
\[A_{\mathrm {L} }^{-1}A=I_{m},\]
where $I_{m}$ is the $m\times m$ identity matrix.
\end{itemize}

Properties:
\bit
\item A matrix doesn't necessarily have a right inverse and doesn't necessarily have a left inverse.
\item A square matrix has a right inverse if and only if it has left inverse if and only if it is invertible.
\item For any invertible matrix, there exists only one right inverse and one left inverse, which are its multiplicative inverse.
\eit
\subsubsection{Generalized inverse or g-inverse (廣義逆)}
A matrix $A^{\mathrm {g} }\in \mathbb {C} ^{n\times m}$ is a generalized inverse of a matrix $A\in \mathbb {C} ^{m\times n}$ if:
\[AA^{\mathrm {g} }A=A.\]

Properties:
\bit
\item For all matrices, there exists at least one generalized inverse, but it's not necessarily unique.
\item For any invertible matrix, there exists only one generalized inverse, which is its multiplicative inverse.
\eit
\subsubsection{Reflexive generalized inverse (自反廣義逆)}
A reflexive generalized inverse of a matrix $\mathbf{A}$ is defined as a matrix $\mathbf{A}^g$ such that:
\[\mathbf{A}\mathbf{A}^g\mathbf{A}=\mathbf{A}\]
\[\mathbf{A}^g\mathbf{A}\mathbf{A}^g=\mathbf{A}^g\]

Properties:
\bit
\item For all matrices, there exists at least one reflexive generalized inverse, but it's not necessarily unique.
\item For any invertible matrix, there exists only one reflexive generalized inverse, which is its multiplicative inverse.
\eit
\subsubsection{Moore–Penrose inverse (摩爾—彭若斯廣義逆) or pseudoinverse (偽逆或擬反)}
A Moore–Penrose inverse of a matrix $\mathbf{A}$ is defined as a matrix $\mathbf{A}^+$ such that:
\[\mathbf{A}\mathbf{A}^+\mathbf{A}=\mathbf{A}\]
\[\mathbf{A}^+\mathbf{A}\mathbf{A}^+=\mathbf{A}^+\]
\[\left(\mathbf{A}\mathbf{A}^{+}\right)^{*}=\mathbf{A}\mathbf{A}^{+}\]
\[\left(\mathbf{A}^+\mathbf{A}\right)^{*}=\mathbf{A}^+\mathbf{A}\]

Properties:
\bit
\item For all matrices, there exists only one Moore–Penrose inverse.
\item For any invertible matrix, there exists only one Moore–Penrose inverse, which is its multiplicative inverse.
\eit
\section{Transformation Matrices}
\subsection{Groups}
\subsubsection{General linear group}
The general linear group of degree $n$ over some field $F$ is the set of $n\times n$ invertible matrices over $F$, together with the operation of matrix multiplication, denoted as $\operatorname{GL}(n,F)$ or $\operatorname{GL}_n(F)$, or simply $\operatorname{GL}(n)$ if the field is understood. It is a topological group and a Lie group under operation $+$.

The Lie algebra of a general linear group is denoted as the symbol of the general linear group with $\operatorname{GL}$ replaced with $\operatorname{\mathfrak{gl}}$.
\subsubsection{Orthogonal group or general orthogonal group}
The orthogonal group or general orthogonal group of degree $n$ over some field $F$ is the set of $n\times n$ invertible matrices over $F$ such that the inverse equals the transpose, together with the operation of matrix multiplication, denoted as $\operatorname{O}(n,F)$ or $\operatorname{O}_n(F)$, or simply $\operatorname{O}(n)$ if the field is understood. An element of an orthogonal group is called an orthogonal matrix. It is a topological group and a Lie group under operation $+$.
\subsubsection{Special orthogonal group or rotation group}
The special orthogonal group or rotation group of degree $n$ over some field $F$ is a subgroup of the orthogonal group of degree $n$ over some field $F$ such that for any matrix in it, the determinant is 1, that is, it preserves the orientation, denoted as $\operatorname{SO}(n,F)$ or $\operatorname{SO}_n(F)$, or simply $\operatorname{SO}(n)$ if the field is understood. An element of a special orthogonal group is called an rotational matrix. It is a topological group and a Lie group under operation $+$.

The Lie algebra of a special orthogonal group is denoted as the symbol of the special orthogonal group with $\operatorname{SO}$ replaced with $\operatorname{\mathfrak{so}}$.
\subsection{Transformation matrices in two-dimensional space}
\subsubsection{Degenerate matrices}
Degenerate to 1 dimension ($(a\lor b)\land (d\lor e)\neq 0$):
\[\begin{bmatrix}a & b\\ca & cb\end{bmatrix}\begin{bmatrix}x \\ y\end{bmatrix}=\begin{bmatrix}ax+by \\ c(ax+by)\end{bmatrix}\]
\[\begin{bmatrix}0 & 0\\d & e\end{bmatrix}\begin{bmatrix}x \\ y\end{bmatrix}=\begin{bmatrix}0 \\ dx+ey\end{bmatrix}\]
Degenerate to 0 dimension:
\[\begin{bmatrix}0 & 0\\0 & 0\end{bmatrix}\begin{bmatrix}x \\ y\end{bmatrix}=\begin{bmatrix}0 \\ 0\end{bmatrix}\]
\subsubsection{Stretching (伸縮/縮放/拉伸)}
A stretch is a linear transformation which scales all distances in a particular direction by a constant factor $k$ but does not affect distances in the perpendicular direction. If $k > 1$, it is really a "stretch"; if $k < 1$, it is a "compression"; if $k = 1$, it has no effect.

The matrix associated with a stretch by a factor $h$ along the $x$-axis and a factor $k$ along the $y$-axis is given by:
\[\begin{bmatrix}
h & 0 \\
0 & k
\end{bmatrix}\]
\subsubsection{Shearing (推移)}
A shearing in the $x$-direction dependent on the $y$ coordination by a factor $k$ is given by:
\[\begin{bmatrix}
1 & k \\
0 & 1
\end{bmatrix}\]
A shearing in the $y$-direction dependent on the $x$ coordination by a factor $k$ is given by:
\[\begin{bmatrix}
1 & 0 \\
k & 1
\end{bmatrix}\]
\subsubsection{Rotation (旋轉)}
A rotation by an angle $\theta$ counterclockwise about the origin is given by:
\[\begin{bmatrix}
\cos(\theta) & -\sin(\theta) \\
\sin(\theta) & \cos(\theta)
\end{bmatrix}\]
\subsubsection{Reflection (鏡射/反射)}
A reflection about a line that passes the original and with $(a,b)=\sqrt{a^2+b^2}(\cos\theta,\sin\theta)$ being a vector in the direction of it is given by:
\[\begin{bmatrix}
\cos(2\theta) & \sin(2\theta) \\
\sin(2\theta) & -\cos(2\theta)
\end{bmatrix}=
\frac{1}{a^2+b^2}\begin{bmatrix}
a^2-b^2 & 2ab \\
2ab & b^2-a^2
\end{bmatrix}\]
\subsubsection{Orthogonal projection (正射影)}
An orthogonal projection about a line that passes the original and with $(a,b)=\sqrt{a^2+b^2}(\cos\theta,\sin\theta)$ being a vector in the direction of it is given by:
\[\begin{bmatrix}
\cos^2(\theta) & \cos(\theta)\sin(\theta) \\
\cos(\theta)\sin(\theta) & \sin^2(\theta)
\end{bmatrix}=
\frac{1}{a^2+b^2}\begin{bmatrix}
a^2 & ab \\
ab & b^2
\end{bmatrix}\]
\subsubsection{Exmaple 1}
\textbf{Question:} The cricle $\mathrm{C}: x^2+y^2=1$ is transformed by $\mathbf{A}=\begin{bmatrix}2&0\\0&3\end{bmatrix}$ to obtain a new curve $\mathrm{C}'$. Find the equation of $\mathrm{C}'$ and its area.

\textbf{Answer:} 
\[\begin{bmatrix}x'\\y'\end{bmatrix}=\mathbf{A}\begin{bmatrix}x\\y\end{bmatrix}=\begin{bmatrix}2x\\3y\end{bmatrix}\]
\[\begin{cases}x'=2x\\y'=3y\end{cases}\Rightarrow \begin{cases}x=\frac{x'}{2}\\y=\frac{y'}{3}\end{cases}\]
\[\qty(\frac{x'}{2})^2+\qty(\frac{y'}{3})^2=1\Rightarrow \frac{x'^2}{4}+\frac{y'^2}{9}=1\]
\[\mathrm{C}'\colon\frac{x'^2}{4}+\frac{y'^2}{9}=1\]
The area of $\mathrm{C}'$ is $\pi \cdot 2 \cdot 3 = 6\pi$. Since $\det(\mathbf{A}) = 6$, and the area of $\mathrm{C}$ is $\pi$, the area of $\mathrm{C}'$ equals $\det(\mathbf{A})$ multiplied by the area of $\mathrm{C}$.
\subsubsection{Exmaple 2}
\textbf{Question:} The line $\mathbf{L}: 4x - 3y = 5$ is transformed in the $x$-direction dependent on the $y$ coordination by twice and then in the $y$-direction dependent on the $x$ coordination by negative theee times to obtain a new line $\mathbf{L}'$. Find the equation of $\mathbf{L}'$.

\textbf{Answer:} 
\[\mathbf{A}=\begin{bmatrix}1&2\\-3&1\end{bmatrix}\]
\[\mathbf{A}^{-1}=\frac{1}{7}\begin{bmatrix}1&-2\\3&1\end{bmatrix}\]
\[\begin{bmatrix}x\\y\end{bmatrix}=\mathbf{A}^{-1}\begin{bmatrix}x'\\y'\end{bmatrix}=\begin{bmatrix}\frac{x'-2y'}{7}\\ \frac{3x'+y'}{7}\end{bmatrix}\]
\[\mathbf{L}': 4\frac{x-2y}{7}+3\frac{3x+y}{7}=5\]
\[\mathbf{L}': 5x+11y+35=0\]
\subsubsection{Example 3}
\textbf{Question:} Reflect the point $\mathrm{P}(5, 2)$ over the line $\mathbf{M}: x - y + 1 = 0$ to obtain a new point $\mathrm{Q}$. Find $\mathrm{Q}$.

\textbf{Answer:} 

Shift one unit to the right: $\mathbf{M}$ becomes $\mathbf{M}':x - y = 0$ and $\mathrm{C}$ becomes $\mathrm{C}'(6, 2)$.

The reflection matrix is: 
\[\mathbf{A} = \begin{bmatrix} 0 & 1 \ 1 & 0 \end{bmatrix}\]
\[\mathrm{D}' = \mathbf{A} \mathrm{C}' = (2, 6)\]

Shift one unit to the left: $\mathrm{D}(1, 6)$
\subsubsection{Example 4}
\textbf{Question:} The line $x+y-2=0$ is transformed by $\mathbf{A}=\begin{bmatrix}a&b\\1&2\end{bmatrix}\in\mathbb{R}^{2\times 2}$ to obtain a new line $2x+3y-4=0$. Find $a$ and $b$.

\textbf{Answer:} 
\[\mathbf{A}\begin{bmatrix}x\\ -x+2\end{bmatrix}=\begin{bmatrix}x'\\ \frac{-2x'+4}{3}\end{bmatrix}\]
\[x-2x+4=\frac{-2x'+4}{3}\]
\[x'=\frac{3}{2}x-4\]
\[ax-bx-2b=x'=\frac{3}{2}x-4\]
\[(a, b)=(-\frac{1}{2}, -2)\]
\subsection{Transformation matrices in three-dimensional space}
\subsubsection{Stretching (伸縮/縮放/拉伸)}
The matrix associated with a stretch by a factor $h$ along the $x$-axis, a factor $k$ along the $y$-axis, and a factor $w$ along the $z$-axis is given by:
\[\begin{bmatrix}
h & 0 & 0 \\
0 & k & 0 \\
0 & 0 & w
\end{bmatrix}\]
\subsubsection{Shearing (推移)}
A shearing in the $x$-direction dependent on the $y$ coordination by a factor $k$ is given by:
\[\begin{bmatrix}
1 & k & 0 \\
0 & 1 & 0 \\
0 & 0 & 1
\end{bmatrix}\]
A shearing in the $x$-direction dependent on the $z$ coordination by a factor $k$ is given by:
\[\begin{bmatrix}
1 & 0 & k \\
0 & 1 & 0 \\
0 & 0 & 1
\end{bmatrix}\]
A shearing in the $y$-direction dependent on the $x$ coordination by a factor $k$ is given by:
\[\begin{bmatrix}
1 & 0 & 0 \\
k & 1 & 0 \\
0 & 0 & 1
\end{bmatrix}\]
A shearing in the $y$-direction dependent on the $z$ coordination by a factor $k$ is given by:
\[\begin{bmatrix}
1 & 0 & 0 \\
0 & 1 & k \\
0 & 0 & 1
\end{bmatrix}\]
A shearing in the $z$-direction dependent on the $x$ coordination by a factor $k$ is given by:
\[\begin{bmatrix}
1 & 0 & 0 \\
0 & 1 & 0 \\
k & 0 & 1
\end{bmatrix}\]
A shearing in the $z$-direction dependent on the $y$ coordination by a factor $k$ is given by:
\[\begin{bmatrix}
1 & 0 & 0 \\
0 & 1 & 0 \\
0 & k & 1
\end{bmatrix}\]
\subsubsection{Rotation (旋轉)}
A rotation by an angle $\theta$ counterclockwise about a line that passes the original and with $(a,b,c)$ being a vector in the direction of it is given by:
\[\frac{1}{a^2+b^2+c^2}\begin{bmatrix}
a^2(1-\cos\theta)+\cos\theta & ab(1-\cos\theta)-c\sin\theta & ac(1-\cos\theta )+b\sin\theta\\
bc(1-\cos\theta )+c\sin\theta & b^2(1-\cos\theta )+\cos\theta & bc(1-\cos\theta )-a\sin\theta\\
ac(1-\cos\theta )-b\sin\theta & bc(1-\cos\theta )+a\sin\theta & c^2(1-\cos\theta )+\cos\theta \end{bmatrix}\]
\subsubsection{Rodrigues' rotation formula}
If $\mb{v},\mb{k}\in\bbR^3$ and $\mb{k}$ is a unit vector describing an axis of rotation about which $\mb{v}$ rotates by an angle $\theta$ according to the right hand rule, the Rodrigues formula for the rotated vector $\mb{v}_r$ is 
\[\mb{v}_r=\mb{v}\cos\theta+(\mb{k}\times\mb{v})\sin\theta+\mb{k}(\mb{k}\cdot\mb{v})(1-\cos\theta).\]
The skew-symmetric matrix $\mb{K}$ in Lie algebra $\operatorname{\mathfrak{so}}(3)$ that generates the one-parameter subgroup of $\operatorname{SO}(3)$:
\[\mb{R}(\theta)=\exp(\theta\mb{K}),\]
such that
\[\mb{v}_r=\mb{R}(\theta)\mb{v},\]
is given by
\[\mb{K}\mb{v}=\mb{k}\times\mb{v},\quad \forall \mb{v}\in\bbR^3,\]
and
\[\mb{R}(\theta)=\mb{I}+\sin\theta\mb{K}+(1-\cos\theta)\mb{K}^2,\]
where $\mb{I}$ is the $3\times 3$ identity matrix.
\subsubsection{Reflection (鏡射/反射)}
A reflection about a plane $ax+by+cz=0$ is given by:
\[\mathbf{I}-\frac{1}{a^2+b^2+c^2}\begin{bmatrix}
2a^2 & 2ab & 2ac \\
2ab & 2b^2 & 2bc \\
2ac & 2bc & 2c^2
\end{bmatrix},\]
where $\mathbf{I}$ is the $3\times 3$ identity matrix.
\subsubsection{Orthogona projection (正射影)}
An orthogonal projection about a plane $ax+by+cz=0$ is given by:
\[\mathbf{I}-\frac{1}{a^2 + b^2 + c^2}\begin{bmatrix}
a^2 & ab & ac \\
ab & b^2 & bc \\
ac & bc & c^2
\end{bmatrix}\]
where $\mathbf{I}$ is the $3\times 3$ identity matrix.
\end{document}

